% VLDB-Template enhances the ACM master template, acmart.cls v2.19 (2026/06/27):
% https://www.acm.org/publications/proceedings-template
% The ACM Latex guide provides further information about the ACM template
%
% This bundle pins acmart.cls v2.19 (rather than relying on whatever version
% ships with an author's local TeX install or Overleaf project) to guarantee
% consistent formatting across all PVLDB submissions.

\documentclass[sigconf, nonacm]{acmart}

\usepackage{algorithm}
\usepackage{algorithmic}
\usepackage{multirow}
\usepackage[normalem]{ulem}
\usepackage{makecell}
\usepackage{comment}
\usepackage[caption=false, font=footnotesize]{subfig}

%%% do not modify the following VLDB block %%
%%% VLDB block start %%%
\usepackage{pvldb}% VLDB/PVLDB formatting overrides, see pvldb.sty for details.
                      % Do not remove -- required for consistent proceedings
                      % formatting across all VLDB submissions.
%%% VLDB block end %%%

%% The following must be set for each paper (values below are placeholders).
%% Everything else (volume/issue/year defaults, the PVLDB Reference Format
%% block, license footer, and Artifact Availability block) is handled by
%% pvldb.sty -- see that file if you need to change shared defaults.
\renewcommand\vldbdoi{XX.XX/XXX.XX}
\renewcommand\vldbpages{XXX-XXX}
% leave empty if no availability url should be set
\renewcommand\vldbavailabilityurl{URL_TO_YOUR_ARTIFACTS} 

\newcommand{\icde}[1]{{\color{black}#1}}
\newcommand{\zj}[1]{{\color{black}#1}}
\newcommand{\done}[1]{{\color{black}#1}}
\newcommand{\kdd}[1]{{\color{black}#1}}
\newcommand{\que}[1]{{\color{black}#1}}
\newcommand{\rem}[1]{{\color{black}#1}}
\newcommand{\rl}[1]{{\color{black}#1}}
\newcommand{\sy}[1]{{\color{black}#1}}
\newcommand{\sydes}[1]{{\color{black}#1}}
\newcommand{\rldes}[1]{{\color{black}#1}}
\newcommand{\rln}[1]{{\color{black}#1}}
\newcommand{\sydesi}[1]{{\color{black}#1}}
\newcommand{\rltwo}[1]{{\color{black}#1}}
\newcommand{\sytwo}[1]{{\color{black}#1}}
\newcommand{\rlend}[1]{{\color{black}#1}}
\newcommand{\syend}[1]{{\color{black}#1}}
\newcommand{\rlx}[1]{\color{black}#1}
\newcommand{\syx}[1]{\color{black}#1}
\newcommand{\syadd}[1]{{\color{black}#1}}
\newcommand{\syexp}[1]{{\color{black}#1}}
\newcommand{\syfix}[1]{{\color{black}#1}}
\newcommand{\syne}[1]{{\color{black}#1}}
\newcommand{\vldb}[1]{{\color{red}#1}}

\newcommand*{\ourmodel}{CAIM}

\begin{document}
\title{Towards Efficient and Interpretable Structured Data Analytics \\ via Causal-Aware Hierarchical Modeling
% Efficient Causal-Aware Hierarchical Modeling for Structured Data
}

%%
%% The "author" command and its associated commands are used to define the authors and their affiliations.

\author{Siyu Chen}
\affiliation{
  \institution{Beijing Institute of Technology}
  \city{Beijing}
  \country{China}
}
\email{3120250945@bit.edu.cn}

\author{Runlin Ma}
\affiliation{
  \institution{Beijing Institute of Technology}
  \city{Beijing}
  \country{China}
}
\email{3220241567@bit.edu.cn}

\author{Zhongle Xie}
\affiliation{
  \institution{Zhejiang University}
  \city{Hangzhou}
  \country{China}
}
\email{xiezl@zju.edu.cn}

\author{Zhaojing Luo}
\affiliation{
  \institution{Beijing Institute of Technology}
  \city{Beijing}
  \country{China}
}
\email{zjluo@bit.edu.cn}

\author{Meihui Zhang}
\affiliation{
  \institution{Beijing Institute of Technology}
  \city{Beijing}
  \country{China}
}
\email{meihui\_zhang@bit.edu.cn}

\author{Xiaokui Xiao}
\affiliation{
  \institution{National University of Singapore}
  \city{Singapore}
  \country{Singapore}
}
\email{xkxiao@nus.edu.sg}

%%
%% The abstract is a short summary of the work to be presented in the
%% article.
\begin{abstract}
\vldb{Relational databases are fundamental data management infrastructures for many real-world applications, including healthcare, finance, and recommendation systems. \vldb{The analytical tasks over structured data stored in relational databases require exploiting complex relationships among columns to improve the quality of downstream data-driven decisions} due to the high-dimensional and sparse nature of the structured data. Existing methods for structured data analytics mainly fall into two categories. The first category utilizes specialized model structures, attention-based or graph-based, to capture inter-column correlations, but these data-driven associations have been shown to include \vldb{spurious dependencies}. The second category incorporates causal analysis into predictive modeling for more reliable relationships, but \vldb{typically} relies on simple \vldb{neural network} architectures like MLPs, limiting their effectiveness in \vldb{modeling complex inter-column dependencies}.

To address these \vldb{challenges}, we propose an efficient Causal-Aware \vldb{hIerarchical} Modeling (CAIM) framework, bridging the gap between high-capacity neural network models and causal analysis. CAIM consists of causal-attentive representation, causal dependency extraction, and hierarchy-aware prediction modules. The causal-attentive representation module \vldb{introduces} a learnable causal matrix to guide the column-wise representation learning. The causal dependency extraction module incorporates a self-supervised relation learning task, regularized by causal structural constraints (i.e., acyclicity and sparsity), to accurately extract causal relationships from complex inter-column interactions. Finally, each column’s representations are aggregated via causal hierarchy-aware weights to perform prediction. \vldb{To improve the efficiency of database-scale analytics, we introduce a lazy-update mechanism that reduces the frequency of costly causal dependency extraction.} Experimental results on eight \vldb{relational} datasets demonstrate that CAIM achieves superior predictive performance \vldb{over} twelve state-of-the-art baselines, \vldb{advancing the capability of database systems to support accurate, interpretable, and scalable data-driven applications.}}
\end{abstract}

\maketitle

%%% do not modify the following VLDB block %%
%%% VLDB block start %%%
\vldbtopmatter
%%% VLDB block end %%%

\section{Introduction}
\rlend{
Structured data, stored in Relational Database Management Systems (RDBMSs)\cite{r:rdbmsnew1, r:rdbmsnew2, r:rdbmsnew3, r:rdbmsnew4, r:rdbmsnew5, r:rdbmsnew6}, 
is widely used in real-world applications, e.g., healthcare \cite{r:healthcarenew1, r:healthcarenew2, r:healthcarenew3}, finance \cite{r:financenew1, r:financenew2, r:financenew3, r:financenew4}, and recommendation \cite{r:recommendationnew1, r:recommendationnew2, r:recommendationnew3, r:recommendationnew4, r:recommendationnew5, r:recommendationnew6}.
}
While Deep Neural Networks (DNNs) have shown great success in unstructured domains, 
their performance often degrades when applied to structured data.
Unlike images or texts, structured data is characterized by high dimensionality and sparsity, which calls for additional pattern learning
\cite{r:patternlearningnew1, r:patternlearningnew2, r:patternlearningnew3}, i.e., mining the relationships between columns \cite{r:relationshipnew1, r:relationshipnew2} that signal the interactions among features. 

\begin{figure}[t]
  \centering
  \includegraphics[width=0.85\linewidth]{causalintr.png}
  \caption{An illustration of the extraction of causal relationships and their application in prediction.}
  \label{fig:causal-righttop}
\end{figure}

Current approaches to learn inter-column relationships within neural models for structured data can be classified into two categories. 
The first leverages specialized structures such as attention \cite{r:tab, r:SAINT} or graph-based architectures \cite{r:gnnbased} to extract patterns between columns. 
However, these data-driven models 
are shown to include spurious patterns\cite{r:superior}, 
\syfix{which are unstable and domain-sensitive
% , and non-transferable 
associations that vary across data distributions and cannot generalize to unseen test data.
% , and improve training performance at the cost of test-time robustness.
}
To solve this, the second category integrates causal analysis 
by jointly learning a causal graph as an auxiliary task \cite{r:CASTLE}. 
These methods \syne{incorporate} regularization \cite{r:linear} to select features that are causally related to the target. 
Since causal relationships are more stable and reflective of true inter-column interactions \cite{r:causalgood, r:pearl2009} than data-driven associations, these approaches often achieve better generalization \syadd{and interpretability}. 
However, these methods are based on simple linear or MLP-based architectures, and their effectiveness in incorporating causal information remains limited\cite{r:linear, r:MLP}.
Consequently, there are still several key challenges in effectively leveraging inter-column relationships for structured data \vldb{analytics}\cite{r:attreg}.

\textbf{Challenge 1:} 
% Modeling and reliably extracting stable inter-column relationships from structured data remains unsolved. The most effective models for structured data \vldb{analytics} today primarily focus on data-driven associations, tending to indiscriminately absorb all statistical correlations between input features and labels in the training data \cite{r:ft-transformer, r:tab, r:SAINT},
\vldb{Automated discovery and management of stable inter-column relationships and dependencies in structured data remain an open problem. 
Existing techniques (e.g., data-driven methods, %functional dependency discovery, 
correlation analysis) primarily focus on data-driven associations, 
% capture statistical associations
}
which may include spurious links that act as shortcuts to improve training performance but worsen test results. 
As illustrated on the left side of Figure~\ref{fig:causal-righttop}, in the diamond price prediction task, \rem{the length, width, and height ($f_1$)} influence the price indirectly through carat ($f_4$). When dataset bias exists, these models may bypass carat and directly associate $f_1$ with the price (e.g., \rem{larger width implies higher price, though a larger width does not necessarily mean a higher carat).} 
Such associations are unstable and fail to generalize across datasets. Causal analysis has demonstrated that causal relationships are more robust and invariant across different settings \cite{r:causalgood}.
However, accurately modeling and extracting causal structures in structured data remains highly challenging, as it requires disentangling complex inter-column relationships and identifying which dependencies are truly important for \vldb{analytics \cite{notears, r:suna}.} \done{Furthermore, the high dimensionality of structured data exacerbates the difficulty of efficiently capturing causal relationships.
}

\textbf{Challenge 2:} 
% There is currently a lack of an effective framework for integrating causal relationship modeling into \rem{high-capacity structured data models}.
\vldb{Integrating discovered causal relationships and dependencies into % **downstream data analytics pipelines** 
downstream data analysis
% — such that attribute representations reflect true causal structure rather than surface correlations — 
remains underexplored 
% challenge 
in structured data analytics and management.
While existing models have attempted to integrate causal knowledge into structured data \vldb{analytics} tasks} \cite{r:CASTLE, r:linear}, they are often built upon simple architectures such as linear models or MLPs. 
These simple additive models, which lack explicit mechanisms for capturing inter-column relationships, have limited capacity to learn complete causal structures\cite{r:MLP}.
To overcome these limitations, a better solution is to incorporate causal analysis into more expressive models, such as attention-based or graph-based architectures, that explicitly model inter-column interactions.
However, \icde{
the structures of these expressive models are different from those of causal modeling, making it challenging to effectively incorporate causal relationships.
}

\textbf{Challenge 3:}
\rem{
There is a lack of an effective approach to exploit the hierarchical causal relationships within structured data.
In our context, the inter-column relationships are characterized through a hierarchical causal graph structure, yet existing methods did not effectively exploit it.
}
As shown on the right side of Figure~\ref{fig:causal-righttop}, 
given the hidden representations of features $E = \{e_1, e_2, \dots, e_n\}$ and an inter-column relational structure \rem{
(As illustrated in the middle of Figure~\ref{fig:causal-righttop}, where each node represents a column in the structured data, the hierarchical causal graph shows that green nodes denote the two-hop causal ancestors of $Y$ and blue nodes represent its direct causal parents.)}, 
existing approaches either treat all feature embeddings equally by concatenating or averaging them \cite{r:tab}, or utilize only partial representations, such as those from the direct causal parents of the target variable \cite{r:CASTLE}.
Both strategies ignore the hierarchical importance among ancestor features, leading to 
degraded predictive performance.
A more reasonable approach is to model the hierarchical importance of features based on their positions in the causal DAG structure, where nodes closer to the target \done{(blue nodes)} typically contribute more to the prediction\vldb{~\cite{r:dagsum}.}
% However, finding an effective way to \zj{incorporating the hierarchical structure of the \rem{structured data causal graph} into the predictive model, so as to fully exploit it to differentiate the importance of features for \rem{structured data prediction}, remains a significant challenge.}
However, finding an effective way to incorporate the hierarchical structure of the structured data causal graph into the predictive model, and thereby fully exploit it to distinguish the importance of features for structured data \vldb{analytics}, remains a significant challenge.

To address these challenges, 
\vldb{
we propose \textbf{C}ausal-\textbf{A}ware H\textbf{i}erarchical \textbf{M}odeling (\ourmodel{}),
}
\rl{a novel \vldb{framework} that bridges the gap between high-capacity neural network models and causal analysis 
for structured data \vldb{analytics}.}
\rl{\ourmodel{} consists of three key components: a causal-attentive representation module, a causal dependency extraction module, and a hierarchy-aware prediction module.}
\sydesi{\rln{First, we design the causal-attentive representation module with a learnable causal matrix that corresponds to a causal graph (i.e., a causal DAG) with columns as nodes and dependencies as edges to guide column-wise representation learning,}}
thereby integrating causal relationship modeling into high-capacity neural network models.
\sydesi{
To optimize the causal matrix, we design the causal dependency extraction module to extract causal structures from the complex inter-column dependencies through a designed self-supervised relation learning task.
}
\rl{
\rln{Last, we design the hierarchy-aware prediction module to aggregate column \rl{representations}, where}
a causal distance-based regularization is 
devised to fully exploit the hierarchical causal DAG structure.}

\rl{To train the causal dependency extraction module with the prediction module properly, \ourmodel{} \rln{is designed with} an alternating optimization strategy. }
\rl{
Moreover, given that the dependency extraction is computationally expensive, we further design a lazy update mechanism that reduces the frequency of updating while maintaining overall model performance.
}
%%% zj-v1: add one to two sentences for interpretation and also in the DAG tech part + also in the contribution

We summarize our contributions as follows:

\begin{itemize}
    \item\que{
    
    % We propose \ourmodel{} \vldb{framework} that integrates causal DAG structures into \sydesi{high-capacity neural architectures} \sy{via a causal-attentive mechanism.} This design enables stable modeling of \done{inter-column} dependencies in structured data for better predictive performances.

    \vldb{We propose \ourmodel{}, a framework that discovers causal-aware \done{inter-column} dependencies and integrates them into downstream structured data analytics via a causal-attentive mechanism.
    This design enables stable modeling and management of inter-column dependencies in structured data for better predictive performance.
    % \ourmodel{} discovers a causal attribute schema from observational data, manages column representations guided by the discovered causal structure, and supports hierarchy-aware analytics that differentiate attribute importance.
    }
    }
    
    \item\que{ 
    We propose an efficient method to enhance causal \sy{DAG learning} from structured data, formulating \sy{causal dependency extraction as an auxiliary task based on a designed self-supervised learning task. }
    Additionally, we \sy{devise causal structural constraints}
    \zj{that ensure valid causal learning and reduce complexity.}}

    \item We design a \sy{hierarchy-aware prediction module} that fully exploits the learned hierarchical causal DAG structure,
    \sy{and the causal distance-based regularization is further proposed to differentiate column importance according to their hierarchical positions in the causal DAG.}

    \item We propose the efficient lazy update strategy to reduce the computational \sy{cost of \ourmodel{} \syadd{by 35\%} via altering the update frequency of causal dependency extraction.}
    
    \item We conduct extensive experiments on \sy{eight} real-world datasets. The results demonstrate that \ourmodel{} outperforms existing SOTA structured data prediction models in terms of \sy{effectiveness, efficiency, and interpretability}.
\end{itemize}

\begin{figure*}[t]
    \centering
    \includegraphics[width=0.92\textwidth]{over1.png}
    \caption{An overview of \ourmodel{}.}
    \label{fig:overview}
\end{figure*}

\section{Related Work}

\subsection{Feature Interaction Modeling}

Traditional machine learning methods, such as XGBoost, AutoML \cite{r:xgboost, r:automl}, \syadd{and more recently TabM \cite{r:tabm},} treat features as independent, limiting their ability to capture complex dependencies. To address this, attention-based models have been proposed. TabTransformer \cite{r:tab} introduces self-attention for feature dependencies, FT-Transformer \cite{r:ft-transformer} improves efficiency through feature tokenization, SAINT \cite{r:SAINT} incorporates inter-sample attention\syadd{, and Orion-Bix \cite{r:orionbix} leverages bidirectional cross‑attention for hierarchical interactions.}
Meanwhile, graph-based methods such as TabGNN \cite{r:gnnbased} and ATT-Reg \cite{r:attreg} explicitly model inter-column relationships via graph structures. 
However, these models are mainly data-driven and prone to capturing spurious dependencies.
%%spurious correlation的定义，原文是Here, a spurious correlation refers to a non-causal statistical association between two variables that exhibits dependence in observational data but lacks an underlying causal mechanism.
%%% (ok) zj-done: (what is this?) remove the explanation here, add in the introduction

\subsection{\vldb{Dependency} and Causal Analysis Integration}
Recent studies introduce \vldb{dependency and} causal reasoning for more reliable \vldb{analytics}\cite{r:causalgraphdb,r:causumx,r:fdx}.
\vldb{
CORDS~\cite{r:cords} discovers statistical correlations between column pairs. HyFD~\cite{r:hyfd} and TANE~\cite{r:tane} mine functional dependencies from relational tables. NOTEARS~\cite{notears} learns causal DAGs via continuous optimization. While these methods can uncover data-driven or causal relationships, they are designed as offline tools and are not designed to be integrated into downstream analytics models, limiting their effectiveness in end-to-end structured data predictive analytics.
}
%%% zj-v-done: 每人讲一句怎么弄
%%% zj-v-done: while they discover data-driven or causal relationships…they are not designed to integrate.. and thus limiting their effectiveness in structured data predictive analytics.
In the meanwhile, LogCause \cite{r:linear} extracts causal relationships via statistical methods and incorporates them as regularization, but relies on a linear prediction model. 
CASTLE \cite{r:CASTLE} jointly learns causal graphs and predictions end-to-end, but \vldb{integrates} into simple MLP backbones. \syfix{
%%% (ok) zj2-done: no need to mention NOTEARS then
LLM‑CD \cite{r:llmcd} combines LLM reasoning with a basic causal discovery algorithm based on conditional independence tests, 
%the PC algorithm,
but its \syne{simple} two‑stage pipeline relies on \vldb{simple downstream models} that lack the capacity to capture complex causal dependencies.
%pipeline prevents joint optimization of causal structure and downstream tasks.
%PC:通过统计条件独立性检验，逐步判断变量之间是否存在边以及边的方向，最终输出一个表示因果关系的有向无环图（DAG）。
} 
Existing causal methods are limited by simple \vldb{downstream prediction model architectures}, restricting their ability to extract complex causal relationships.
\vldb{Our proposed \ourmodel{} integrates causal modeling with the expressive capacity of high-capacity models,} 
enabling stable extraction of causal inter-column relationships for robust \vldb{predictive analytics}.

\section{Problem Formulation}

\sy{
Given a table $T = \{(x_1, x_2, \dots, x_n, y)_i \mid i = 1, 2, \dots, I\}$, where each of the $I$ samples consists of $n$ feature columns and a label $y$. \syfix{Each column can be either a numerical feature or a categorical feature and 
we thus also use the term "feature" to refer to "column". 
The} goal is to train a predictive function $f: \mathbf{x} \to y$. 
This objective aims to minimize the error between the predicted $\hat{y}$ and the ground truth $y$, requiring the model to effectively capture the complex patterns inherent in high-dimensional and sparse structured data.}

\section{Methodology}

\vldb{
In modern structured data analytics tasks, analysts often need to understand how attributes (columns) depend on one another. 
% For example, in healthcare database analytics, distinguishing whether a medication causally improves an outcome versus merely correlating with it is critical. 
Traditional database techniques (e.g., correlation analysis, functional dependency mining) and recent high-capacity neural network models capture statistical associations between pairwise columns yet these data-driven approaches often include spurious correlations \syfix{which are unstable and vary across data distributions and} hinder generalization.
% capture statistical associations but cannot distinguish causal from spurious relationships. This section presents our framework for automated causal dependency profiling and causality-guided structured data analytics.
}
\syadd{
In the \syne{meantime}, causal analysis uncovers \syfix{robust} and invariant relationships between features and targets\cite{r:causalitydb}. 
% Unlike purely correlational patterns, causal features 
\syfix{\syne{This} offers a promising solution: integrating causal relationships
% , they maintain robust and consistent associations with the prediction target independent of data distribution shifts. This enables
with the data-driven model to focus on features with direct and 
stable effects. 
% improving generalization and reducing misleading patterns.
}
}
\syfix{However, to achieve this, how to extract reliable causal relationships from observational structured data 
% is challenging due to spurious correlations and hidden confounders. Moreover,
and how to integrate them
% the extracted causal graph 
into the end-to-end differentiable model with a completely different structure 
% while preserving structural constraints 
\syne{pose} great challenges. 
% To our knowledge, no existing framework fully addresses both challenges in a unified manner.
}

% Our main idea is that existing methods fail to integrate causal analysis with high-capacity models, resulting in suboptimal performance in both causal discovery and prediction. 
% To address this limitation, 
% we propose \ourmodel{}, which 
% Thus, our main idea is to combine an expressive neural architecture with causal learning for robust structured data prediction.
% to bridge causal structural reasoning for structured data prediction.

First, 
regarding causal relationship extraction, 
as causal relationships 
\icde{
describe how each feature is causally generated \syne{by} other features as its causal parents,
} 
we design the self-supervised reconstruction task (Section~\ref{sec:causal_extraction})
\syadd{
which reveals the underlying generative mechanisms among columns, as located in the upper-right part of Figure~\ref{fig:overview},
} to learn \syfix{the causal structure} by predicting each column from the others. 
% \syadd{
% To extract sparse causal links and determine causal directions, this reconstruction task randomly masks certain columns and predicts them from others,
% } combined with causal structural constraints, specifically acyclicity and sparsity, ensuring valid and sparse causal graph learning.
Second, regarding causal relationship integration, 
\syfix{
since causal relationships exist between features
} 
and the self-attention mechanisms also capture inter-column relationships, we thus design the causal-attentive representation module (Section~\ref{sec:causal_attentive}) to integrate causal relationships with the self-attention mechanism, which is shown in the center of Figure~\ref{fig:overview}. 
% shows the \ourmodel{} architecture and as shown in the center of the figure} 
% \syfix{, this module fuses the causal mask with self-attention to capture causally meaningful inter-column relationships.}
% this module directly processes the column embeddings from the left through a dual-path Causal-Aware Attention Block,} \syadd{where the attention scores and the causal mask are fused via element-wise multiplication to form causal attention,} rectifying spurious correlations and ensuring that learned representations reflect true causal structure.
%% It is worth noting that our main focus is not necessarily strict causal discovery itself, but rather exploring how causal structure can be used to support predictive modeling in downstream structured data analytics tasks. In this sense, the learned causal graph is regarded as an auxiliary component that may help improve the prediction task.
Third, since causal structures are inherently hierarchical and thus features contribute unequally to prediction, we design the hierarchy-aware prediction module (Section~\ref{sec:causal_hierarchy_prediction}) \syadd{located in the lower-right part of Figure~\ref{fig:overview}, whose key idea is \syne{to} 
% \syadd{to exploit this hierarchical property \syfix{via a} causal distance-based regularization, 
align each column's predictive importance with its causal position and encourage causally closer features to exert greater influence while still allowing distant features to contribute.}
% , this module operates alternately with the reconstruction task.
\syfix{
It is worth noting that our primary objective \syne{for} \ourmodel{} is not strict causal discovery, but rather the exploratory use of the learned causal structure as auxiliary knowledge to approximately identify causally important features, thereby improving predictive performance on downstream structured data analytics tasks.
}

% \syadd{Before delving into the detailed design of each component, we first summarize the core mathematical symbols used throughout this section in Table~\ref{tab:symbols}.}
% \vldb{
% \begin{table}[t]
%     \centering
%     \caption{TABLE OF SYMBOLS}
%     \label{tab:symbols}
%     \begin{tabular}{c|l}
%         \hline
%         \textbf{Symbol} & \textbf{Description} \\ \hline
%         $S = [x_1,\dots,x_n,y]$ & Input vector with $n$ features and target $y$ \\ \hline
%         $E = [e_1,\dots,e_{n+1}]$ & Embedding matrix after field embedding \\ \hline
%         $d_e$ & Embedding dimension \\ \hline
%         $C \in \mathbb{R}^{(n+1)\times(n+1)}$ & Learnable causal matrix \\ \hline
%         $M = \text{Sigmoid}(C)$ & Soft causal mask (sigmoid of $C$) \\ \hline
%         $T$ & Thresholded adjacency matrix of DAG \\ \hline
%         $d_i$ & Causal distance of feature $i$ to target $y$ \\ \hline
%         $D_{\max}$ & Large constant for non-ancestor nodes \\ \hline
%         $n_h$ & Number of attention heads \\ \hline
%         $Q, K, V$ & Query, key, value matrices in self-attention \\ \hline
%         $\text{Attention}_h$ & Attention scores for head $h$ \\ \hline
%         $\text{FixedAttention}_h$ & Attention with causal mask applied \\ \hline
%         $H \in \mathbb{R}^{(n+1)\times d_e}$ & Causal-attentive hidden representation \\ \hline
%         $N$ & Number of attention blocks \\ \hline
%         $\hat{s}_i$ & Reconstructed feature value \\ \hline
%         $\mathcal{L}_{\text{rec}}$ & Reconstruction loss \\ \hline
%         $\mathcal{L}_{\text{DAG}}$ & Acyclicity regularization loss \\ \hline
%         $\mathcal{L}_{\text{sparse}}$ & Sparsity regularization loss \\ \hline
%         $\alpha$ & Reconstruction loss weight \\ \hline
%         $\beta$ & Structural constraint strength \\ \hline
%         $\mathcal{L}_R$ & Total loss for causal extraction module \\ \hline
%         $a_i$ & Learnable aggregation weight for column $i$ \\ \hline
%         $h_{\text{pred}}$ & Aggregated representation for prediction \\ \hline
%         $\hat{y}$ & Predicted target \\ \hline
%         $\mathcal{L}_{\text{pred}}$ & Prediction loss \\ \hline
%         $\lambda$ & Weight decay strength \\ \hline
%         $\mathcal{L}_{\text{dis}}$ & Distance-aware regularization \\ \hline
%         $\mathcal{L}_P$ & Total loss for prediction module \\ \hline
%     \end{tabular}
% \end{table}
% }
\subsection{Preprocessing Module}
\rlend{
For structured data with $n$ columns, the input is represented as $\mathbf{S} = [x_1, \ldots, x_n, y] \in \mathbb{R}^{(n+1)}$. 
\syfix{
% For a classification or regression task on structured data with $n$ features, this input vector serves as the base representation. 
A \texttt{[msk]} token is used to replace the column to be predicted.}
After that, the resulting vector $\mathbf{S}' \in \mathbb{R}^{n+1}$ is transformed into an embedding matrix $E = [\mathbf{e}_1, \dots, \mathbf{e}_{n+1}]$, where each $\mathbf{e}_i \in \mathbb{R}^{d_e}$. 
Specifically, categorical columns are mapped via embedding lookup, where each column $i$ has an embedding matrix $E_i \in \mathbb{R}^{m_i \times d_e}$ ($m_i$ is the count of category values). 
% \syadd{This ensures that each category in each field has its own unique embedding.}
Numerical columns are transformed via $e_j = x_j \cdot \mathbf{e}^0_j$, where $x_j$ is the standardized \syadd{(zero-mean, unit-variance)} input value and $\mathbf{e}^0_j \in \mathbb{R}^{d_e}$ is a learnable base embedding. 
% \syadd{This design preserves the multiplicative nature of numerical values.}
\syadd{This linear projection design is commonly adopted in structured data deep learning\cite{r:armnet, r:autoint, r:ft-transformer}, as it preserves both the scalar magnitude of numerical values and the field-specific semantics of each column.
}
\syadd{The \texttt{[msk]} token is mapped to a shared embedding $e_{msk}$.}
% Through this process, the input vector is transformed into an embedding matrix $\mathbf{E} = [\mathbf{e}_1, \ldots, \mathbf{e}_{n+1}]$, where each $\mathbf{e}_i \in \mathbb{R}^{d_e}$ represents the embedding of the $i$-th field.
}

\subsection{Causal Analysis Theory}\label{sec:causal_theory}

\sytwo{\ourmodel{}} is based on the assumption that \syadd{the data generation process follows the nonparametric structural equation model (NPSEM), under which there exists a causal DAG} among the input features and the target\cite{r:pearl2009,r:peters2017}. 
% \syadd{This is a well-justified assumption in causal inference\cite{r:pearl2009,r:peters2017}.} 
In this causal framework, features $\mathbf{X}$ and the target $y$ form a DAG $G = (V, E)$, where each $v \in V$ corresponds to a column \sytwo{(including both features and the target)} and edge $\langle u, v \rangle \in E$ indicates that $u$ is a direct causal parent of $v$. \rl{The causal DAG is characterized by acyclicity and \syadd{an explicit sparsity assumption.}} 
The dependencies among \syfix{variables (in our setting, columns)} 
are assumed to follow NPSEM\cite{r:pearl2009}, where each non-root node is generated from its causal parents:
\begin{equation}
    v_i = f_i(\mathrm{Pa}(v_i), \epsilon_i), \quad i \in [n+1]
    \label{eq:structural-eq}
\end{equation}
where $\mathrm{Pa}(v_i)$ denotes the direct causal parents of $v_i$, $\epsilon_i$ is an independent noise term, and $f_i$ can be any nonlinear function.

\rln{
In our design, \syend{a learnable causal matrix parameterizes the DAG and} is \syend{integrated} into the attentive mechanism of the model to rectify spurious correlations between \syne{variables}, \syend{as detailed }in \vldb{the causal-attentive representation module (Section~\ref{sec:causal_attentive}).}
% \icde{It is worth noting that our primary objective is not necessarily strict causal discovery, but rather leveraging causal structure to improve predictive performance on downstream structured data analytics tasks.}
Moreover, the design of the causal dependency extraction module (Section~\sytwo{\ref{sec:causal_extraction}) follows the NPSEM assumption that each variable is generated from its causal parents. \syadd{This is} instantiated via a self-supervised reconstruction task \syadd{that predicts each feature from the others, thereby learning the causal dependencies among \syne{variables}.}
}
}

\subsection{Causal-Attentive Representation Module}\label{sec:causal_attentive}

\rlend{For structured data, attention-based mechanisms serve as the core module in many modern architectures for capturing inter-column dependencies\cite{r:autoint, r:tabm}, but the learned attention scores may include spurious associations.} \rln{To address this, we first devise a learnable causal matrix (Section~\ref{sec:causal_mask_gen}), 
and then integrate it with the self-attention mechanism to rectify spurious associations (Section~\ref{sec:causal_attn_block}). }\syend{ The detailed structure of this module is illustrated in Figure~\ref{fig:overview}B.}

% \subsubsection{Causal Analysis Theory}

% \kdd{In the standard supervised learning setting, we denote the input feature variables by $\mathbf{X} = [X_1, \ldots, X_d] \in \mathcal{X}$ and the target variable by $Y \in \mathcal{Y}$, where $\mathcal{X} \subseteq \mathbb{R}^d$ is a $d$-dimensional feature space and $\mathcal{Y} \subseteq \mathbb{R}$ is a one-dimensional target space. Let $P_{\mathbf{X},Y}$ denote the joint distribution of the features and target. We observe a dataset $\mathcal{D} = \{(\mathbf{X}^{(i)}, Y^{(i)})\}_{i=1}^N$ consisting of $N$ i.i.d. samples drawn from $P_{\mathbf{X},Y}$. The goal of a supervised learning algorithm $\mathcal{A}$ is to find a predictive model $f_Y: \mathcal{X} \rightarrow \mathcal{Y}$ in a hypothesis space $\mathcal{H}$ that can explain the association between the features and the target variable. When the hypothesis space $\mathcal{H}$ is too complex or the training data fails to adequately represent the underlying distribution $P_{\mathbf{X},Y}$, overfitting may occur---manifested as a mismatch between training and testing performance of $f_Y$. This motivates the usage of regularization to reduce the complexity of $\mathcal{H}$ so that $\mathcal{A}$ will learn a function with better generalization. The choice of regularization is typically guided by assumptions about the underlying data structure.}


% \kdd{Our primary goal is therefore not causal discovery, but to enhance prediction performance by learning a graph with causal properties (directed, acyclic, and sparse) through appropriate constraints, which collectively contribute to improved generalization.}


\subsubsection{Causal Mask Generator}\label{sec:causal_mask_gen}

\syadd{To model the causal relationships among \syfix{features}, we design a mask generator with a set of learnable parameters
and a function 
that \syne{maps} to a directed acyclic graph (DAG) % matrix representing 
of causal dependencies. 
Specifically, we define} \rl{a learnable causal matrix} $C \in \mathbb{R}^{(n+1) \times (n+1)}$, 
where each entry $C_{i,j}$ represents the strength of the directed edge from node $j$ to node $i$ in the causal graph.
\done{This directed matrix models the potential causal influences among \syfix{features}, allowing the model to capture both the hierarchical (direct and indirect) dependencies and the directionality of their relationships.
}
To transform the continuous weights in $C$ into a usable causal mask, \syend{we exploit an element-wise sigmoid function, defined as $M = \text{Sigmoid}(C)$, where a higher value indicates a stronger causal effect in that direction.}
\icde{
In the early stage of training, the causal mask is dense, exerting a weak modulating effect on the model, so the model's behavior approximates that of a standard attention network. 
As training progresses, the mask gradually becomes sparse, thereby progressively introducing causal constraints.
}

\subsubsection{\vldb{Causal-Aware} Attention Block
}\label{sec:causal_attn_block}

\sydesi{
Given the designed causal mask, we aim to incorporate its causal structural knowledge into the column learning process. 
To this end, we design the \vldb{Causal-Aware} Attention block.
\rl{First, the embedding matrix $E \in \mathbb{R}^{(n+1) \times d_e}$ is projected into query ($Q$), key ($K$), and value ($V$) matrices, yielding raw attention scores $\text{Attention}_{h} = \text{Softmax}(\frac{Q_h K_h^\top}{\sqrt{d_k}})$.}
}
\sydesi{However, raw attention scores primarily capture statistical associations, which may include spurious links. 
So we integrate the causal DAG by elementwise multiplying the attention weights with the causal mask $M$:}
\begin{equation}
    \text{FixedAttention}_h = \text{Attention}_h \odot M
    \label{eq:masked-attn}
\end{equation}
This multiplicative operation effectively integrates causal priors into the attention mechanism, \sydesi{allowing each column to attend only to its causally relevant predecessors. Consequently, this design emphasizes }causally meaningful dependencies.
\icde{
Note that downstream features can still be highly informative, which motivates our design of multiplicative soft masking rather than hard pruning: multiplication does not fully mask out the original attention but reweights it to enhance causal predecessors, thereby allowing the model to leverage both upstream and downstream information.
% We acknowledge that in a single training distribution, downstream variables can still be highly informative. 
If a direction (whether upstream to downstream or downstream to upstream) has a strong response in either the causal mask or the attention, their product remains large, and thus the information is not discarded.
}

The output of each attention head is computed as:
\begin{equation}
    \text{Attended}_h = \text{FixedAttention}_h \cdot V_h
    \label{eq:attended-head}
\end{equation}
\syadd{The causal mask M is shared across all attention heads and layers. This design choice stems from the fact that the causal graph captures a global, intrinsic data-generating mechanism among \syfix{features}, which is independent of specific heads or layers. Although different heads learn diverse representations through their own projection matrices (Q, K, V), the underlying causal relationships must be respected by all heads. Hence, sharing the mask does not restrict representational diversity; each head can still capture distinct feature interaction patterns while adhering to the same causal constraints.}

\rlend{
The outputs from all heads are then concatenated and passed through an output projection: \syadd{
\begin{equation}
    \textit{Attended} = \text{Concat}_{h=1}^{n_h}(\textit{Attended}_h) W_O
    \label{eq:concat-proj}
\end{equation}}
where $W_O \in \mathbb{R}^{d_e \times d_e}$ is a projection matrix. Finally, the result goes through residual connections and layer normalization, yielding the causal-attentive representation $H \in \mathbb{R}^{(n+1) \times d_e}$.
}

\sy{The transformation from the embedding matrix $E$ to the latent representation is denoted as $\Phi(\cdot)$\syend{, with $H^{(i)}$ representing the output of }the $i$-th representation module. After passing through $N$ such modules, \syadd{where $N$ denotes the number of stacked attention blocks,} we obtain the final causal-attentive representation $H \in \mathbb{R}^{(n+1) \times d_e}$ \syend{via the composition}} \syadd{:
\begin{equation}
    H = \Phi_N \circ \cdots \circ \Phi_1(E)
    \label{eq:final-composition}
\end{equation}}

\subsection{Causal Dependency Extraction Module}\label{sec:causal_extraction}
\rln{According to NPSEM, each \syadd{variable} (including features and a label) is generated from its causal parents via a function $f_i$. 
This insight guides the optimization of the learnable causal matrix $M$: by learning to predict each \syadd{feature} from others, we can discover the causal parent set and thus optimize the causal matrix. Based on this insight, we design the causal dependency extraction module with a reconstruction task. }
\syadd{
Specifically, Equation~\ref{eq:structural-eq} states each variable is generated from its causal parents, and reconstruction directly mirrors this by predicting each column from the others. 
Beyond reconstruction, we further introduce causal structural constraints (acyclicity and sparsity) to ensure the learned graph better aligns with the NPSEM assumptions.
} 
To this end, a random column is \sytwo{replaced with a $[msk]$ token to generate a masked vector $S'$, which is mapped to an embedding matrix $E$ where the $[msk]$ token is represented by a shared learnable embedding $\mathbf{e}_{msk}$. The} task is to predict the masked input $s_i$ using the aggregated representation $H_i$ from the remaining columns.

\syfix{
Recall that Equation~(\ref{eq:attended-head}) defines the per-head output as a matrix product.
For each target position $i$, the $i$-th row of this output is a weighted sum over all source positions $j$:
\begin{equation}
    \text{Attended}_h[i] = \sum_{j=1}^{n+1} \text{FixedAttention}_h[i,j] \cdot V_h[j]
    \label{eq:causal-decomposition}
\end{equation}
where $\text{FixedAttention}_h[i,j]$ is the causal-masked attention weight from column $j$ to column $i$ in head $h$, which reflects the strength of the edge $\langle j, i \rangle$ in the learned causal DAG.
After multi-head concatenation and a residual connection, we denote the final causal-attentive representation of column $i$ as $H_i$.} 

\rl{To further ensure a valid causal structure, we introduce causal structural constraints (i.e., acyclicity and sparsity), detailed in the following sections}\syadd{, and the overall structure is illustrated in Figure~\ref{fig:subfig1}.}

\begin{figure}[h]
    \centering
    \includegraphics[width=\columnwidth]{causal dependency extraction module.png}
    \caption{Overview of the Causal Dependency Extraction Module.}
    \label{fig:subfig1}
\end{figure}

\subsubsection{Reconstruction Task}
Given the masked column's representation $H_i$, \syend{the value} $s_i$ is reconstructed via a feed-forward network \syend{as $\hat{s}_i = \text{FFN}(H_i)$.} 
\syfix{The intuition is that this reconstruction task implicitly learns the NPSEM structural equations $v_i = f_i(\mathrm{Pa}(v_i), \epsilon_i)$ by predicting each variable from the others. Minimizing the reconstruction error under sparsity and acyclicity constraints makes the causal DAG identifiable, thus implementing the NPSEM formulation.}
The final reconstruction loss is the Cross-Entropy or MSE between the predicted value $\hat{s}_i$ and the ground truth $s_i$, depending on whether the column is categorical or numerical. To balance the scale discrepancy between categorical and numerical variables, 
we introduce a weighting $\eta$ to ensure fair training:\syne{
\begin{equation}
    \mathcal{L}_{\text{rec}} =
    \begin{cases}
    \eta \cdot \|\hat{s}_i - s_i\|_2^2, & \text{if } \delta_i = 0 \\
    -\log \hat{s}_i[s_i], & \text{if } \delta_i = 1
    \end{cases}
    \label{eq:reconstruction-loss}
\end{equation}}
where $\delta_i=1$ indicates a categorical variable and $\delta_i=0$ indicates a numerical variable.

\subsubsection{Causal Structural Constraints}
\rl{To enforce that the learned causal matrix satisfies the causal DAG structure assumed in NPSEM, we introduce two structural constraints: an \sy{acyclicity loss} ensuring valid directed acyclic structure, and a sparsity loss encouraging sparse causal parent identification.}

\paragraph{Acyclicity(DAG) Loss}
The \sy{acyclicity loss} penalizes the cycle in the Causal Mask $M$, which is designed as:

\begin{equation}
    \mathcal{L}_{\text{DAG}} = \mathrm{tr}(e^M) - (n+1)
    \label{eq:dag-loss}
\end{equation}
Since the matrix exponential can be expanded as a power series: $e^M = I + M + \frac{1}{2!}M^2 + \cdots$, the trace term $\mathrm{tr}(e^M)$ can then be rewritten as $\mathrm{tr}(I) + \mathrm{tr}(M) + \mathrm{tr}(\frac{1}{2!}M^2) + \cdots$, where $\mathrm{tr}(M^k)$ indicates the number of cycles of length $k$ in the graph. 
Therefore, any cycle in the graph will contribute a positive value to $\mathcal{L}_{\text{DAG}}$, making it an effective regularizer to enforce the DAG constraint.
% \icde{
% Since $M=\mathrm{Sigmoid}(C)$ has entries within $[0,1]$ and is pushed towards zero by the sparsity loss, numerical overflow of $e^M$ is not observed in practice. Moreover, vanishing gradients are expected when $M$ is nearly acyclic, as the DAG constraint is largely satisfied, consistent with the convergence behavior that gradients naturally diminish near a valid DAG.
% }

\paragraph{Sparsity Loss}
To promote a sparse causal graph and encourage clear identification of root-level features (i.e., variables with few or no parents), we add a group-Lasso penalty on the rows of \syne{the soft causal mask $M$} (which correspond to in-degrees):
\begin{equation}
    \mathcal{L}_{\text{sparse}} = \sum_{i=1}^{n+1} \|M_{i, :}\|_2
    \label{eq:sparsity-loss}
\end{equation}
This structured sparsity helps drive some rows entirely to zero, effectively selecting a sparse set of causal edges and promoting the detection of root nodes in the DAG.

\subsubsection{Overall Loss for Causal Dependency Extraction Module}
The total loss of the \sy{causal dependency extraction module} is:
\begin{equation}
    \mathcal{L_{R}} = \alpha\mathcal{L}_{\text{rec}} + \beta(\mathcal{L}_{\text{DAG}} + \mathcal{L}_{\text{sparse}})
    \label{eq:total_reconstruction_loss}
\end{equation}
where $\alpha$ and $\beta$ are hyperparameters designed to balance the reconstruction task with the regularization terms, as well as to coordinate the importance between the \sy{causal dependency extraction module} and the prediction \kdd{module} introduced later.

\subsection{Hierarchy-Aware Prediction Module}
\label{sec:causal_hierarchy_prediction}

\rlend{
Through the causal dependency extraction module, we learn the causal DAG defining the hierarchical relationships among variables, and obtain hidden representations $H$ for each column. 
Traditional aggregation methods (e.g., concatenation or averaging) typically ignore this hierarchical structure and \syne{do} not fully utilize the extracted hierarchical causal DAG knowledge. 
% \syadd{To fully exploit this causal knowledge to mine significant structural patterns for downstream prediction,} 
To this end, we design a hierarchy-aware weighting mechanism that bridges structural reasoning and feature fusion.} \syadd{The overall structure is illustrated in Figure~\ref{fig:subfig2}.}

\begin{figure}[h]
    \centering
    \includegraphics[width=\columnwidth]{hierarchy-aware prediction module.png}
    \caption{Overview of the Hierarchy-Aware Prediction Module.}
    \label{fig:subfig2}
\end{figure}


\subsubsection{Predictor Architecture}
\rlend{
The predictor aggregates each row of the input $H \in \mathbb{R}^{(n+1) \times d_e}$ into a global representation $h_{\text{pred}}$ using learnable weights $\{a_1, a_2, \dots, a_{n+1}\}$. 
Note that $H$ is computed with the target $y$ replaced by the $[msk]$ token, and its embedding $\mathbf{e}_{n+1}$ is the shared learnable embedding $\mathbf{e}_{msk}$. The global representation is calculated as:
\begin{equation}
    h_{\text{pred}} = \sum_{i=1}^{n+1} a_i h_i
    \label{eq:h_pred}
\end{equation}
}


$h_{\text{pred}}$ is fed into a feedforward neural network to produce the prediction: $\hat{y} = \text{FNN}(h_{\text{pred}})$.
The prediction loss $\mathcal{L}_{\text{pred}}$ depends on task type (regression or classification), and is defined as:
\begin{equation}
    \mathcal{L}_{\text{pred}} =
    \begin{cases}
    \| \hat{y} - y \|_2^2, & \text{for regression tasks} \\
    -\log \hat{y}[y], & \text{for classification tasks}
    \end{cases}
    \label{eq:pred_loss}
\end{equation}

\syfix{
The weight $a_i$ in Equation~\ref{eq:h_pred} is crucial as it measures the importance of column $i$. 
We argue that this importance is inherently related to the causal distance from each feature to the target $y$. 
In the following, we first compute the distance $d_i$ from the learned DAG, and then design a distance-aware regularization that down-weights features with larger $d_i$.
}

\subsubsection{Causal Information Extraction}
\syadd{
Empirically, we binarize the learned causal matrix $M$ into the adjacency matrix $T$ by applying a threshold of $0.5$
%: $T_{i,j} = 1$ if $M_{i,j} > 0.5$, otherwise $0$. 
This ensures that only sufficiently confident causal relationships are retained for hierarchy-aware prediction.}
Then, to extract the distance from each node to the label node $y$, we leverage matrix exponentiation. Specifically, the $k$-th power $T^k[:, n+1]$ encodes the number of paths of length $k$ from any node to $y$. If $T^k[i, n+1] > 0$, it implies that node $i$ can reach $y$ in exactly $k$ steps. Using this property, we construct a distance vector $d \in \mathbb{R}^{n+1}$, where each entry $d_i$ represents the path length from node $i$ to $y$ in the DAG:
\begin{equation}
    d_i = 
    \begin{cases}
    \min\{k \mid T^k[i, n+1] > 0\}, & \text{if such } k \text{ exists} \\
    D_{\max}, & \text{otherwise}
    \end{cases}
    \label{eq:distance_matrix}
\end{equation}
\noindent
Here, $D_{\max}$ is a large constant assigned to nodes that are not ancestors of $y$ (e.g., descendants or unrelated nodes). % \syadd{Under the NPSEM assumption, direct causes are the primary drivers of the target, and the shortest-path distance naturally captures this causal hierarchy.} 
This distance-aware representation enables subsequent modules to focus more on nodes that are causally closer to the label.

\subsubsection{Distance-aware Regularization}
\label{sec:reg}

\icde{
A feature's causal distance to the target reflects its importance: distance 1 corresponds to direct causes, and larger distances to indirect ones. 
As importance naturally decays with distance\cite{r:pearl2009,r:peters2017}, which is consistent with prior work prioritizing direct parents\cite{r:CASTLE,r:linear}, 
the distance thus captures this hierarchical structure, i.e., features causally closer to the target are more prediction-relevant.
}
Consequently, given the computed distance vector $d \in \mathbb{R}^{n+1}$, we introduce a distance-aware regularization term to \syadd{align the predictor weights with the causal hierarchy.} 
We modify the weight decay strength on the predictor coefficients $\{a_1, \dots, a_{n+1}\}$ by incorporating the causal distances. The regularization term is defined as:
\begin{equation}
    \mathcal{L}_{\text{dis}} = \sum_{i=1}^{n+1} \lambda d_i a_i^2
    \label{eq:distance_regularization}
\end{equation}

\noindent
where $\lambda$ is the original weight decay strength. 
\syadd{
We do not evaluate \syfix{predictive importance} only with causal distance, nor restrict the predictor to use only direct parents. 
Equation~\eqref{eq:distance_regularization} implements a soft regularization: features with larger causal distance are down-weighted but never forced to zero. 
The hyperparameter $\lambda$ controls regularization strength. 
When $\lambda = 0$, all features are treated equally, recovering data-driven aggregation. 
This preserves flexibility to exploit useful distant signals while encouraging prioritization of causally closer features.
}

\subsubsection{Overall Loss for Prediction \kdd{module}}
In the prediction \kdd{module}, the Causal Matrix is frozen and the overall loss is:
\begin{equation}
    \mathcal{L}_{P} = \mathcal{L}_{\text{pred}} + \mathcal{L}_{\text{dis}}
    \label{eq:total_prediction_loss}
\end{equation}

\subsection{Optimization for \ourmodel{}}
\rlend{
For the \ourmodel{} model, there are two \syne{special} sets of parameters to optimize: the causal matrix $C$ (\icde{causal dependency extraction module}) and the predictor weights $a_i$ (\icde{hierarchy-aware prediction module}). 
According to Equations~\ref{eq:distance_matrix} and \ref{eq:distance_regularization}, the optimization of $a_i$ depends on the learned $C$
\syfix{
which defines the soft mask $M$ binarized to obtain the adjacency matrix $T$. 
The distance $d_i$ is then computed from $T$ and controls $a_i$ via regularization.
} 
Therefore, we propose to update these two sets of related parameters in an alternating manner.
}
% \syadd{through the causal dependency extraction module and the hierarchy-aware prediction module, respectively. 
\syadd{
In this section, we detail the optimization steps and further introduce a Lazy Update strategy designed to improve training efficiency.
}

\subsubsection{Alternating Optimization Strategy}
%%% (ok) zj1-done: delete algoritm 1 and only reserve algorithm 2
In each iteration, we first update the causal dependency extraction module using loss $\mathcal{L}_R$ (Equation~\ref{eq:total_reconstruction_loss}), which optimizes the causal matrix $C$ to learn a valid and sparse causal graph. 
Subsequently, we update the prediction module using loss $\mathcal{L}_P$ (Equation~\ref{eq:total_prediction_loss}), which optimizes the predictor weights $a_i$ while keeping $C$ fixed.  
\syadd{This alternating process enables the model to simultaneously improve feature reconstruction and predictive performance, achieving joint optimization. It also disentangles causal structure learning from target prediction, ensuring that the learned causal relations are unaffected by target-specific noise and thus more generalizable.}
\icde{
Note that the updates of both $\mathcal{L}_R$ and $\mathcal{L}_P$ are within a unified alternating framework, optimized via stochastic gradient procedures.
This preserves the convergence property of the overall alternating algorithm. 
% Moreover, the reconstruction objective provides stronger supervision for refining causal dependencies, which further stabilizes the optimization process near convergence.
}
    
\subsubsection{Lazy Update}

\syadd{Updating the causal dependency extraction module can be time-consuming due to the need to compute the \syfix{acyclicity} loss. 
\syfix{
However, we observe that after the first few epochs, the causal matrix $C$ and the associated acyclicity loss change only gradually, so they do not need to be updated at every mini-batch.} 
To reduce the overall computational cost, we thus propose a lazy update strategy where the core idea is to reduce the frequency of updating the costly causal dependency extraction module.} 
% \syfix{By updating the causal dependency extraction module less frequently, we can significantly reduce the computational overhead without sacrificing predictive performance.} 
\syadd{Let $C$ denote the parameters related to the causal mask, \syfix{$a_i$ denote the predictor weights}, 
and $\Theta$ \syne{represent} all other model parameters.
$T_s$ is the number of mini-batches used for updating the causal dependency extraction module, and $T_c$ is the cycle length.} 
For every \syadd{$T_c$} mini-batches, the first $T_s$ are used to update the causal dependency extraction module \vldb{ (lines 8--12)}, while the remaining \syadd{$T_c - T_s$} update the \syadd{hierarchy-aware} prediction module \vldb{(lines 14--18)}. 
%多讲一些动机目的,参考GM-reg
A smaller \syadd{$T_s/T_c$} reduces training time but may lower performance. Empirically, we set $T_s=2$ and \syadd{$T_c=10$}.

\begin{algorithm}[h]
    \caption{Lazy \vldb{Update} of \ourmodel{}}
    \label{alg:lazy_update}
    \small
    \begin{algorithmic}[1]
        \STATE \textbf{Input}: Training data $\mathcal{D}$, causal mask matrix $C$, the predictor weights $a_i$, other parameters $\Theta$
        \STATE \textbf{Parameters}: causal update steps $T_s$, total cycle length $T_c$
        \STATE \textbf{Output}: Trained \ourmodel{} model
        \STATE Initialize iteration counter $t \leftarrow 0$
        \FOR{each mini-batch $B \in \mathcal{D}$}
            \STATE Embed each sample $\mathbf{S}$ into $\mathbf{E}$ as described in the \textit{Preprocessing module} section
            \IF{$t \bmod T_c < T_s$}
                \STATE /*Causal Dependency Extraction*/
                \STATE Randomly select a column in $B$ and replace it with \texttt{[msk]}
                \STATE Compute hidden representations $H$ through Eq.~\ref{eq:masked-attn}--Eq.~\ref{eq:final-composition}
                \STATE Compute reconstruction objective $\mathcal{L}_R \leftarrow \alpha \mathcal{L}_{\text{rec}} + \beta(\mathcal{L}_{\text{DAG}} + \mathcal{L}_{\text{sparse}})$
                \STATE Update $C$ and $\Theta$ by minimizing $\mathcal{L}_R$
            \ELSE
                \STATE /*Hierarchy-aware Prediction*/
                \STATE Replace target $y$ in $B$ with \texttt{[msk]}
                \STATE Compute hidden representations $H$ through Eq.~\ref{eq:masked-attn}--Eq.~\ref{eq:final-composition}
                \STATE Fix causal mask $C$ and compute prediction objective $\mathcal{L}_P \leftarrow \mathcal{L}_{\text{pred}} + \mathcal{L}_{\text{dis}}$
                \STATE Update $a_i$ and $\Theta$ by minimizing $\mathcal{L}_P$
            \ENDIF
            \STATE $t \leftarrow t + 1$
        \ENDFOR
        \STATE \textbf{return} Trained \ourmodel{} model
    \end{algorithmic}
\end{algorithm}

% Model complexity is an important consideration for practical deployment. To facilitate analysis and reduce the influence of hyperparameter variations, we denote the embedding dimension as $d_e$, the number of attention heads as $K$, the number of input fields as $m$, and the total parameters of all feed-forward networks as $n_f$.  

% For the embedding layer, categorical features are represented by embedding tables with a parameter size of $O(Md_e)$, where $M$ denotes the total number of categories. Numerical features are projected through linear transformations with complexity $O(md_e)$. \rem{Since $m/M$ is usually very small, the latter can be safely ignored. Furthermore, the forward propagation complexity of the embedding layer is only $O(m)$, which is negligible in comparison to that of the subsequent modules.}
 

% In the attention module, each head requires projection matrices of size $O(d_e^2)$, leading to a total parameter size of $O(Kd_e^2)$. The training complexity is $O(Km^2d_e)$, while inference can be reduced to $O(Kmd_e)$ through pre-computation.  

% The prediction module is mainly implemented by feed-forward networks, whose total parameter size is summarized as $n_f$, and both training and inference complexity scale proportionally with $n_f$.  

% Additional structures, such as relation-aware masks or structural constraints, typically introduce only $O(m^2)$ parameters. For example, when computing pairwise distances or DAG-based regularization, the complexity in principle could reach $O(km^2)$, where $k$ is the step size of dependency exploration. In practice, however, only the top-$k$ relations (with $k$ fixed) are considered, so the effective complexity reduces to $O(m^2)$. Moreover, lazy update strategies can further reduce the overhead of acyclicity loss during training. Importantly, since $M$ is fixed, these operations impose no additional cost during inference.  

% \done{
% Overall, the dominant parameter scale is $O(Md_e + Kd_e^2 + n_f)$. Since $d_e$ is typically fixed (e.g., 1024), the parameter scale grows linearly with the number of categories $M$. The training complexity is approximately $O(Km^2d_e)$, increasing quadratically with the number of input fields $m$, while the inference complexity is $O(Kmd_e)$, linear in $m$. These complexities are comparable to those of standard Transformer-based models. 
% Moreover, $m$ is usually small (often below $10^3$), so the overall computation cost remains modest.
% }

% The lazy update strategy is pivotal in translating the aforementioned theoretical efficiency into practical training acceleration. This strategy mitigates the computational overhead of the reconstruction branch by updating it only once after every ten updates of the prediction branch (i.e., a 1:10 ratio). This deliberate scheduling effectively amortizes the dominant O(m²) cost of computing the acyclicity loss across multiple batches. Consequently, the strategy significantly reduces the average training time per iteration without compromising convergence or final performance. 

\begin{table*}
    \caption{Performance Comparison on Real-World Applications}
    \label{tab:model_comparison}
    \centering
    \resizebox{\textwidth}{!}{
    \begin{tabular}{l l c c c c c c c c c c c c c}
      \toprule
      Dataset & Metric & XGBoost & AutoML & \syadd{TabM} & FT-Trans & Tab-Trans & SAINT & \syadd{Orion-Bix} & \sydesi{ATT-Reg} & TabGNN & CASTLE & LogCause & \syfix{LLM-CD} & \ourmodel{} \\
      \midrule
      CreditCard
        & AUC (\%)
        & \syne{97.88$\pm$0.04}
        & \syne{97.81$\pm$0.04}
        & \syne{98.10$\pm$0.04}
        & \syne{98.17$\pm$0.03}
        & \syfix{97.48$\pm$0.04}
        & \syne{97.98$\pm$0.03}
        & \syne{97.78$\pm$0.05}
        & \syfix{97.27$\pm$0.04}
        & \syne{98.14$\pm$0.03}
        & 95.80$\pm$0.08
        & 96.48$\pm$0.05
        & \syfix{97.26$\pm$0.05}
        & \textbf{98.62$\pm$0.02} \\
      Cardiovascular
        & AUC (\%)
        & \syne{78.96$\pm$0.01}
        & \syne{78.90$\pm$0.01}
        & \syne{79.05$\pm$0.02}
        & \syne{79.08$\pm$0.02}
        & \syne{79.04$\pm$0.02}
        & \syne{79.08$\pm$0.02}
        & \syne{78.94$\pm$0.03}
        & \sydesi{75.10$\pm$0.02}
        & \sydesi{77.21$\pm$0.02}
        & \syfix{78.35$\pm$0.04}
        & 77.62$\pm$0.03
        & \syne{78.61$\pm$0.03}
        & \syadd{\textbf{79.37$\pm$0.02}} \\
      Adult
        & AUC (\%)
        & \syfix{90.57$\pm$0.04}
        & \syfix{90.55$\pm$0.06}
        & \syne{90.93$\pm$0.05}
        & \syfix{90.97$\pm$0.04}
        & \syfix{90.47$\pm$0.06}
        & \syfix{90.94$\pm$0.05}
        & \syfix{90.49$\pm$0.06}
        & \syfix{90.25$\pm$0.05}
        & \sydesi{89.50$\pm$0.05}
        & 88.55$\pm$0.08
        & 87.87$\pm$0.04
        & \syfix{86.18$\pm$0.06}
        & \syadd{\textbf{91.16$\pm$0.04}} \\
      \midrule
      Diamonds
        & RMSE
        & \syne{0.147$\pm$0.01}
        & \syne{0.146$\pm$0.01}
        & \syne{0.144$\pm$0.01}
        & \syne{0.142$\pm$0.01}
        & 0.156$\pm$0.01
        & \syne{0.143$\pm$0.01}
        & \syne{0.151$\pm$0.01}
        & \sydesi{0.171$\pm$0.01}
        & \syne{0.147$\pm$0.01}
        & 0.184$\pm$0.03
        & 0.197$\pm$0.02
        & \syfix{0.194$\pm$0.02}
        & \syadd{\textbf{0.137$\pm$0.01}} \\
      Elevator
        & RMSE
        & \syfix{0.739$\pm$0.02}
        & \syfix{0.737$\pm$0.02}
        & \syfix{0.734$\pm$0.02}
        & \syfix{0.733$\pm$0.02}
        & \syfix{0.734$\pm$0.02}
        & \syfix{0.731$\pm$0.02}
        & \syadd{0.742$\pm$0.02}
        & \sydesi{0.753$\pm$0.02}
        & \sydesi{0.736$\pm$0.02}
        & 0.908$\pm$0.03
        & 0.916$\pm$0.04
        & \syfix{0.904$\pm$0.03}
        & \syadd{\textbf{0.722$\pm$0.02}} \\
      HouseSale
        & RMSE
        & \syne{0.173$\pm$0.01}
        & \syne{0.174$\pm$0.01}
        & \syne{0.173$\pm$0.01}
        & 0.172$\pm$0.01
        & 0.171$\pm$0.01
        & 0.173$\pm$0.01
        & \syadd{0.177$\pm$0.01}
        & \sydesi{0.184$\pm$0.01}
        & \sydesi{0.176$\pm$0.01}
        & 0.189$\pm$0.02
        & 0.195$\pm$0.01
        & \syfix{0.180$\pm$0.02}
        & \textbf{0.165$\pm$0.01} \\
      Crime
        & RMSE
        & \sydesi{0.588$\pm$0.02}
        & \sydesi{0.586$\pm$0.02}
        & \syadd{0.566$\pm$0.02}
        & \sydesi{0.565$\pm$0.01}
        & \sydesi{0.570$\pm$0.02}
        & \sydesi{0.560$\pm$0.01}
        & \syadd{0.585$\pm$0.02}
        & \sydesi{0.572$\pm$0.02}
        & \sydesi{0.576$\pm$0.02}
        & \sydesi{0.762$\pm$0.02}
        & \sydesi{0.951$\pm$0.03}
        & \syfix{0.759$\pm$0.02}
        & \syadd{\sydesi{\textbf{0.557$\pm$0.01}}} \\
      MEPS
        & RMSE
        & \sydesi{0.920$\pm$0.02}
        & \sydesi{0.911$\pm$0.02}
        & \syadd{0.903$\pm$0.02}
        & \sydesi{0.904$\pm$0.02}
        & \sydesi{0.906$\pm$0.02}
        & \sydesi{0.897$\pm$0.01}
        & \syadd{0.914$\pm$0.02}
        & \sydesi{0.907$\pm$0.02}
        & \sydesi{0.912$\pm$0.02}
        & \sydesi{0.922$\pm$0.02}
        & \sydesi{1.024$\pm$0.03}
        & \syfix{0.956$\pm$0.03}
        & \syadd{\sydesi{\textbf{0.889$\pm$0.01}}} \\
      \bottomrule
    \end{tabular}
    }
\end{table*}

\section{Experiments}

\subsection{Baselines}
\syfix{\ourmodel{} integrates causal structural knowledge into attention-based deep learning and machine learning architectures, and we thus include machine learning~\vldb{(e.g., XGBoost\cite{r:xgboost}, AutoML\cite{r:automl}, and TabM\cite{r:tabm})}, attention-based~\vldb{(e.g., FT-Transformer\cite{r:ft-transformer}, TabTransformer\cite{r:tab}, SAINT\cite{r:SAINT}, and Orion-Bix\cite{r:orionbix})} and causal-based~\vldb{(e.g., CASTLE\cite{r:CASTLE}, LogCause\cite{r:linear}, and LLM-CD\cite{r:llmcd})} methods as baselines.
% to capture stable inter-column relationships.
Additionally, \ourmodel{} exploits the learned hierarchical causal graph to differentiate feature importance. 
We thus also include graph-based~\vldb{(e.g., ATT-Reg\cite{r:attreg} and TabGNN\cite{r:gnnbased})} methods as baselines.}
% Given this multi-faceted design,} we compare \ourmodel{} with \syfix{twelve} baselines of four categories \syadd{, including machine learning, attention-based, graph-based, and causal-based methods.} 

%%% zj-v-done: merge the above two paragraphs

\vldb{
\begin{itemize}
    \item \textbf{XGBoost}: A gradient boosting decision tree model that builds an ensemble of weak learners through sequential optimization, known for its high efficiency and strong performance on structured data.
    
    \item \textbf{AutoML}: An automated machine learning framework that performs model selection and hyperparameter optimization automatically, aiming to achieve competitive performance without manual tuning.
    
    \item \textbf{TabM}: A modern structured data deep learning model that incorporates advanced regularization and architecture designs for structured data prediction.
    
    \item \textbf{FT-Transformer}: A transformer-based model for structured data prediction that uses column-wise feature tokenization and row-wise attention for contextual learning.
    
    \item \textbf{TabTransformer}: A model that applies transformers to encode categorical features via column attention.
    
    \item \textbf{SAINT}: A self-attention-based model that introduces both row-wise and column-wise attention to better capture complex feature interactions in structured data.
    
    \item \textbf{Orion-Bix}: A model that leverages bidirectional cross-attention and hierarchical feature interaction learning for structured data.
    
    \item \textbf{ATT-Reg}: A graph-based method that models pairwise inter-column relationships via attentive weight matrices.
    
    \item \textbf{TabGNN}: A graph-based method that constructs a multiplex graph to model sample relationships.
    
    \item \textbf{CASTLE}: A causal-based method that learns the causal directed acyclic graph (DAG) as an adjacency matrix embedded in the input layer of an MLP network.
    
    \item \textbf{LogCause}: A causal-based method that regularizes a linear predictor with the causal strengths between columns, computed using statistical methods.
    
    \item \textbf{LLM-CD}: A causal-based method that leverages LLM-derived causal priors for causal graph discovery, yet relies on a simple prediction model.
\end{itemize}
}

\subsection{Datasets}
We select \sydesi{eight} real-world datasets for evaluation, including three classification and \sydesi{five} regression tasks. \syadd{The selected datasets span small, medium, and large numbers of features.}
% \sydesi{Notably, we incorporate the Crime and MEPS datasets to simulate high-dimensional scenarios, allowing us to assess \ourmodel{}’s effectiveness in environments with a large number of features.}
% All datasets are publicly available on Kaggle\footnote{\url{https://www.kaggle.com/datasets}} or the UCI Machine Learning Repository\footnote{\url{https://archive.ics.uci.edu/datasets/}}.

\noindent \textbf{CreditCard}: \vldb{This dataset is sourced from a bank transaction database that records credit card transactions. It contains 284,807 samples with 30 numerical features for binary classification of fraudulent transaction.}

\noindent \textbf{Cardiovascular}: \vldb{This dataset is sourced from a clinical database that collects patient health records. It contains 70k samples with 11 categorical features for binary classification of heart disease.}

\noindent \textbf{Adult}: \vldb{This dataset is sourced from the U.S. Census Bureau database, representing a typical public-sector database table. It contains 48,842 samples with 8 categorical and 6 numerical features for binary classification of whether income exceeds \$50K.}

\noindent \textbf{Diamonds}: \vldb{This dataset is sourced from an online diamond listings platform that records diamond transactions. It contains 53,940 samples with 3 categorical and 6 numerical features for regression of diamond price.}

\noindent \textbf{Elevator}: \vldb{This dataset is sourced from an industrial sensor database that monitors elevator operations. It contains 112,001 samples with 7 numerical features for regression of elevator vibration.}

\noindent \textbf{HouseSale}: \vldb{This dataset is sourced from online real estate listing platforms that records housing sales in Indian cities. It contains 30,138 samples with 40 categorical features for regression of house sale price.}

\noindent \textbf{Crime}: \vldb{This dataset is sourced from a combined socio-economic, law enforcement, and crime statistics dataset that compiles community crime statistics. It contains 1,994 samples with 122 numerical features for regression of violent crime rate.}

\noindent \textbf{MEPS}: \vldb{This dataset is sourced from a national healthcare survey database that tracks medical expenditures and patient demographics. It contains 48,982 samples with 138 mixed demographic and medical features for regression of healthcare utilization.}

\syadd{\subsection{Hyperparameter Settings}
For a fair comparison, all deep learning methods share \syfix{the} same training configurations. Specifically, the hidden size is set to $64$, the batch size is set to $128$, and the learning rate is set to $1\times10^{-4}$. In all the experiments, the optimizer is Adam with a weight decay of $1\times10^{-5}$, and the number of layers is set to two for all the deep learning models. The training epochs for all the datasets and methods are set to $100$. The dropout rate is set to $0.1$, 
\syfix{and for transformer-based methods, the per-head dimension is set to the hidden size divided by 4\cite{r:attention}.}
}

\rlend{\subsection{Configurations and Metrics}}

All methods are implemented using PyTorch 2.2.2 with CUDA 12.2, and the experiments are conducted on a server equipped with an Intel(R) Xeon(R) w7-3455 CPU and two NVIDIA RTX A6000 GPUs.
\syend{We employ AUROC and RMSE to evaluate the model performance on classification and regression tasks, respectively. }

% In all experiments, the batch size is set to 128. The optimizer \sy{employed} is Adam with a weight decay of $1 \times 10^{-5}$. Other hyperparameters are selected based on the best-performing configuration for each model on each dataset.

% \subsection{Evaluation metrics}
% We \done{\sy{employ}} AUROC and RMSE to evaluate the model performance on classification and regression tasks, respectively. 
% A higher AUROC indicates better classification performance, while a lower RMSE implies better regression performance.

\subsection{Comparison on Real World Datasets}

\rln{In this section, we compare \ourmodel{} with \syfix{twelve} baselines across eight real-world datasets. }
\icde{Every experiment is repeated 5 times with random seeds, and we report mean performance.}
As summarized in Table~\ref{tab:model_comparison}, \ourmodel{} consistently outperforms all baselines across both classification and regression tasks.
%%\done{Specifically, the two causal-based models perform the worst due to their \done{exploit} of relatively simple backbone networks, such as MLPs or linear predictors, which restrict their ability to capture complex nonlinear dependencies among features and thus limit their overall representational power.

\rln{
Traditional machine learning methods (XGBoost and AutoML) perform reasonably on low-dimensional numerical datasets such as CreditCard and Elevator, but show limitations on high-dimensional or mixed-feature datasets (e.g., Crime, MEPS), as they do not explicitly model \sydesi{inter-column relationships}.
\syne{TabM modestly improves over tree-based methods but lags behind attention-based models, underscoring the need for explicit inter-column modeling.} 
\sytwo{Transformer-based methods (FT-Transformer, Tab-Transformer, \syfix{Orion-Bix,} and SAINT) achieve the second-best performance overall. 
\syfix{
Orion-Bix with bi-axial cross-attention shows no clear advantage on these benchmarks because its bidirectional interactions introduce redundant feature dependencies.
} 
FT-Transformer, Tab-Transformer and SAINT perform competitively on both classification and regression tasks, benefiting from their strong expressive capacity for modeling inter-column interactions. 
However, these top-performing methods rely purely on data-driven correlations without explicit causal constraints, which limits further improvement.
Graph-based methods (ATT-Reg and TabGNN) explicitly model feature relationships but show \sydesi{limited effectiveness} across most datasets, indicating that \sydesi{data-driven associations are} insufficient without proper structural constraints.
\syfix{Causal-based methods with simplistic designs (CASTLE, LogCause \syfix{and LLM-CD}) perform the worst among all baselines, } particularly on high-dimensional datasets such as MEPS.
%检查一下这一章和对应表格
We attribute this to their reliance on simple \sydesi{architectures} (MLPs or linear predictors), which \syne{restrict} their ability to capture complex nonlinear dependencies.}
% On Adult, \syadd{its advantage over the second best (91.30\% vs. 91.15\%) is statistically significant, confirming that even small gains on saturated datasets are reliable. 

\icde{
In contrast, \ourmodel{} achieves the best results because it integrates causal reasoning with high‑capacity attention modeling.
}
The causal-attentive representation module rectifies spurious associations learned by attention, while the hierarchy-aware prediction module exploits the causal hierarchy to \sydesi{differentiate importance of features}, leading to superior performance.
}

\subsection{Interpretability}
In this section, we demonstrate \ourmodel{}’s interpretability through two representative scenarios. 
% \syadd{We do not claim that the learned DAGs are guaranteed to recover the true causal graph under strict identifiability conditions on real-world data. Within our framework, the causal graph serves as an auxiliary structure to aid prediction; even when the necessary conditions for causal identifiability are not fully met, the impact on the primary prediction task remains limited.} 
Specifically, we examine whether the learned causal DAGs on the Adult and Diamonds datasets align with domain knowledge about feature importance. We analyze predictive paths in the DAGs and show the results in Figure~\ref{fig:interpretability}.
\rlend{
On the left side (Adult), the learned DAG shows that \textit{age}, \textit{education-num}, and \textit{relationship} have the strongest influence on income. Moreover, \ourmodel{} captures secondary causal patterns such as \textit{education} influencing \textit{education-num}, and \textit{marital-status} and \textit{sex} determining \textit{relationship}. 
This hierarchy is sociologically intuitive: higher education leads to more schooling years, which increases earning potential\cite{r:hout2012}.
On the right side (Diamonds), \textit{clarity} and \textit{carat} emerge as the primary direct causes of \textit{price}, while \textit{length}, \textit{width}, and \textit{height} indirectly influence price through \textit{carat}\cite{r:diamonds}. 
\ourmodel{} avoids learning spurious relationships such as \textit{sex} directly determining \textit{income} or \textit{length} directly determining \textit{price}.}
This demonstrates that \ourmodel{} can uncover causal structures that are robust, interpretable, and consistent with expert understanding of real-world generative mechanisms.
\icde{
Since there is no numerical groundtruth for causal relationships in the real-world applications like sociology and gemology, we primarily focus on qualitative visualization to present the learned causal structures and compare them with known conclusions from the literature\cite{r:hout2012, r:diamonds}. 
For quantitative evaluation, we will conduct synthetic data experiments in the next section.
}


\subsection{Synthetic Data Experiments}
To evaluate the causal discovery capability of \ourmodel{}, we conduct experiments on synthetic datasets
with known causal DAGs~\sydesi{and compare it against causal baselines, i.e., CASTLE\cite{r:CASTLE}, LogCause\cite{r:linear}} \syexp{and NOTEARS\cite{notears}.
% , a classic gradient-based causal discovery method.
}
\syexp{We consider two settings: linear Gaussian (consistent with the NPSEM assumption) and nonlinear Gaussian (more realistic, violating the linearity assumption).}
We first present the linear experiments in detail, followed by the nonlinear experiments.

\subsubsection{Linear Synthetic Data Experiments}

\rltwo{
We generate synthetic datasets following the NPSEM framework\syend{\cite{r:pearl2009, r:CASTLE}}.
Specifically, we first construct a random DAG structure with $n$ nodes representing $n-1$ features and label $y$, where edges denote causal dependencies. Each edge is assigned a random weight sampled from a uniform distribution. Following the topological order of the DAG, we initialize root nodes with random values sampled from $\mathcal{N}(0, 1)$, and then generate each non-root node from its causal parents according to the structural equation $v_i = \sum_{j \in \mathrm{Pa}(v_i)} w_{ji} v_j + \epsilon_i$, where $\epsilon_i \sim \mathcal{N}(0, \sigma^2)$ is independent noise. We generate datasets with 5-node (6 edges) and 10-node (15 edges) DAG structures, and train \ourmodel{} and the baselines on these datasets to assess their ability to accurately recover the underlying causal graph.
}

\begin{figure}[t]
    \centering
    \subfloat[Adult]{%
        \includegraphics[width=0.45\linewidth,height=0.8in]{inter1.png}%
        \label{fig:inter1}%
        % \vspace{-3mm}%
    }
    \hspace{0.02\linewidth}
    \subfloat[Diamonds]{%
        \includegraphics[width=0.45\linewidth,height=0.8in]{inter2.png}%
        \label{fig:inter2}%
        % \vspace{-3mm}%
    }
    \caption{Interpretability analysis results on two datasets.}
    \label{fig:interpretability}
\end{figure}

\rltwo{
We adopt three evaluation metrics: True Positive Rate (TPR), the ratio of correctly identified to all true edges; False Positive Rate (FPR), the ratio of falsely inferred to all non-edges; and RMSE, the root mean squared error for downstream task prediction.
}

%%% zj-v-done: tables 2-4 can now be smaller? like tables 5 & 6?
\begin{table}[h]
  \caption{Performance Comparison of Causal Structure Discovery on Linear Synthetic Datasets}
  % \vspace{-3mm}
  \label{tab:synthetic}
  \centering
  \small
  \setlength{\tabcolsep}{3pt}
  \begin{tabular}{l ccc ccc}
    \toprule
    Model
    & \multicolumn{3}{c}{5-Node DAG}
    & \multicolumn{3}{c}{10-Node DAG} \\
    \cmidrule(lr){2-4} \cmidrule(lr){5-7}
      & TPR $\uparrow$ & FPR $\downarrow$ & RMSE $\downarrow$
      & TPR $\uparrow$ & FPR $\downarrow$ & RMSE $\downarrow$ \\
    \midrule
    LogCause
      & 0.67 & 0.43 & 0.2312
      & 0.53 & 0.31 & 0.2156 \\
    CASTLE
      & 0.67 & 0.36 & 0.2059
      & 0.60 & 0.24 & 0.1949 \\
    NOTEARS
      & \textbf{1.00} & \textbf{0.00} & --
      & \textbf{1.00} & \textbf{0.00} & -- \\
    \textbf{\ourmodel{} (Ours)}
      & 0.83 & 0.21 & \textbf{0.1678}
      & 0.80 & 0.19 & \textbf{0.1871} \\
    \bottomrule
  \end{tabular}
\end{table}
% \vspace{-3mm}

\syfix{As shown in Table~\ref{tab:synthetic}, NOTEARS achieves near-perfect causal graph recovery (TPR=1.0, FPR=0.0), which is expected since the linear Gaussian data are generated under the NPSEM assumptions for which NOTEARS is designed. 
As a pure causal discovery method, however, NOTEARS cannot perform downstream prediction with RMSE marked as ``--''. 
More importantly, even a perfectly recovered causal graph does not directly translate to optimal prediction: as we will show in Section~\ref{sec:end_to_end}, feeding NOTEARS' output as the causal mask into \ourmodel{}'s architecture yields substantially worse results than \ourmodel{}'s end-to-end joint learning of causal relationships and downstream task prediction.
%%% zj1-done: the reason "This is because ..." moved to H? or delete it
Among the methods that jointly perform causal discovery and prediction, i.e., CASTLE and LogCause, \ourmodel{} consistently achieves higher TPR and lower FPR, indicating more accurate causal discovery, while achieving better RMSE in downstream prediction.
}
\sydesi{The simple architectures of CASTLE and LogCause limit their}
ability to capture indirect dependencies, leading to false or missing edges.
In contrast, \ourmodel{}’s attentive causal masking integrates causal reasoning with the attention mechanism, enabling causal relations to build on inter-column interactions and reinforce one another.
\sydesi{
% Overall, experimental results validate that \ourmodel{} consistently outperforms all baselines in both causal structure discovery and
In terms of downstream prediction
% . By achieving higher TPR and lower FPR while maintaining superior accuracy,
\ourmodel{} also achieves superior performance than baselines, demonstrating its hierarchy-aware design that effectively identifies the dominant causal structure necessary for reliable prediction.}

\subsubsection{Nonlinear Synthetic Data Experiments}

\syexp{To better reflect real-world scenarios, we further generate a second set of synthetic datasets with \syfix{10-node (15 edges) }DAGs under a nonlinear Gaussian setting.
The graph structure and noise distribution remain the same, but each non-root node is generated by a nonlinear function of its parents:
$v_i = f_i(Pa(v_i)) + \epsilon_i$, where $\epsilon_i \sim \mathcal{N}(0, \sigma^2)$. Here $f_i$ is a randomly initialized MLP with one hidden layer (32 units, tanh activation). All other parameters (noise variance, graph sparsity, etc.) are kept identical to the linear case.
% Results for this nonlinear experiment are shown in Table~\ref{tab:nonlinear}.
}

\syadd{
\begin{table}[h]
  \caption{Performance Comparison of Causal Structure Discovery on Nonlinear Synthetic Datasets}
  % \vspace{-3mm}
  \label{tab:nonlinear}
  \centering
  \small
  \setlength{\tabcolsep}{4pt}
  \begin{tabular}{l ccc}
    \toprule
    Model & TPR $\uparrow$ & FPR $\downarrow$ & RMSE $\downarrow$ \\
    \midrule
     LogCause   & 0.47 & 0.39 & 0.2512 \\
     CASTLE     & \syne{0.50} & 0.28 & 0.2186 \\
     NOTEARS   & 0.53 & 0.23 & -- \\
     \textbf{\ourmodel{} (Ours)} & \textbf{0.73} & \textbf{0.20} & \textbf{0.1836} \\
    \bottomrule
  \end{tabular}
\end{table}
}
%% \vspace{-3mm}

\syfix{As shown in Table~\ref{tab:nonlinear}, all methods degrade under the nonlinear setting compared to their linear counterparts, since the nonlinear parent-child relationships are more challenging to recover. Notably, NOTEARS, which achieved perfect recovery under linear Gaussian data (TPR=1.0, FPR=0.0), suffers the most significant degradation, as its linear SEM assumption no longer holds under the tanh-MLP nonlinear generation process. CASTLE and LogCause also exhibit performance drops, constrained by their simple backbone architectures. In contrast, \ourmodel{} maintains the best TPR and FPR among all methods, demonstrating that its causal-attentive masking generalizes beyond the linear NPSEM assumption to more complex and real nonlinear data generation processes.}


\subsection{Effectiveness of End-to-end Learning}\label{sec:end_to_end}
Integrating causal structure into predictive modeling \syne{can follow} a two-stage approach---first discovering a causal graph via a dedicated algorithm and then feeding it as input to the predictor---or an end-to-end approach that jointly learns the causal structure and the prediction task. \ourmodel{} adopts the latter. To evaluate whether optimization provides benefits, we construct a two-stage baseline by applying NOTEARS\cite{notears}, a gradient-based causal discovery method, to the Adult and Diamonds datasets. The recovered graphs are binarized and used as causal masks $M$ in \ourmodel{}'s architecture. 
In contrast, the full \ourmodel{} learns $M$ jointly with the prediction objective.

\begin{table}[h]
  \caption{Performance Comparison Using NOTEARS Graph as Causal Mask}
  % \vspace{-3mm}
  \label{tab:notears_mask}
  \centering
  \small
  \setlength{\tabcolsep}{2pt}
  \begin{tabular}{lcc}
    \toprule
    Variant & Adult (AUC $\uparrow$) & Diamonds (RMSE $\downarrow$) \\
    \midrule
    NOTEARS (as causal mask) & 0.8797 & 0.149 \\
    \textbf{\ourmodel{} (Ours)} & \syne{\textbf{0.9116}} & \syne{\textbf{0.137}} \\
    \bottomrule
  \end{tabular}
\end{table}
% \vspace{-3mm}

As shown in Table~\ref{tab:notears_mask}, the NOTEARS variant substantially underperforms \ourmodel{} on both datasets. We attribute this gap to two factors. \syfix{First, NOTEARS is designed for causal discovery, not prediction. 
As a result, it learns sparse and binary-like causal graphs, which may remove weak edges that contribute useful information for prediction.}
Second, learning the causal structure \icde{separately from the downstream task} 
yields a static mask that cannot adapt to the prediction signal during training. 
In contrast, \ourmodel{}'s end-to-end learning enables mutual reinforcement between causal structure and predictive representations:
the reconstruction task guides causal discovery, while the prediction loss retains task-relevant edges that pure causal discovery might suppress.
\icde{
% Admittedly, NOTEARS may recover causal graphs more accurately, yet its downstream performance is poor. 
This confirms \ourmodel{}'s core philosophy of exploiting causal knowledge as an auxiliary for downstream prediction rather than pursuing exact causal discovery, which leads to better predictive performance.
}

\subsection{Distribution Shift Experiments}

\kdd{\sydesi{Since causal modeling is motivated by robustness under distribution shifts},
we conduct out-of-distribution experiments on the Diamonds dataset to evaluate whether \ourmodel{} achieves this goal. We consider two representative settings: a feature-level shift and a spurious-correlation reversal scenario inspired by Simpson's Paradox\cite{r:peters2017}.

\rln{
For the feature-level shift, we split the training and test sets by the Color feature: diamonds with colors D, E, or F constitute 90\% of the training set, while those with colors G, H, I, or J constitute 90\% of the test set. 
\rlx{
For the spurious-correlation reversal scenario, we construct a setting where the correlation between a feature and the target reverses between training and test due to a confounding variable. 
Specifically, in the Diamonds dataset, both Carat and Clarity have positive causal effects on Price (the target). 
However, we introduce a confounding pattern by biased sampling: the training set is dominated by samples where Carat and Clarity are negatively associated (small-carat with high-clarity, and large-carat with low-clarity), while the test set is dominated by samples where they are positively associated. 
Since Carat is a stronger predictor for Price, this creates a spurious negative correlation between Clarity and Price in the training set (Pearson $r$ = -0.42).
% , which reverses to positive in the test set ($r$ = +0.44). 
A model that relies on this spurious correlation will fail under this distribution shift. 
}
}

\rlend{
We compare \ourmodel{} with \ourmodel{} w/o Causal, which removes the causal modeling components while keeping the same backbone architecture. \ourmodel{} consistently outperforms \ourmodel{} w/o Causal in both settings: for the feature-level shift, \ourmodel{} achieves RMSE of 0.198 compared to 0.204 (\ourmodel{} w/o); for the Simpson's Paradox setting, \ourmodel{} achieves 0.223 compared to 0.244 (\ourmodel{} w/o). 
These results demonstrate that causal modeling improves robustness under feature distribution shift and avoids spurious correlations}\syfix{, as the acyclic sparse DAG captures causal relations instead of surface associations.} 
\rlx{Notably, \ourmodel{} identifies the true positive causal effect from Clarity to Price, without being misled by the induced spurious negative correlation in the training set.} 
%These results demonstrate that \ourmodel{} is more robust to distribution shifts. While the non-causal variant exploits unstable statistical associations, \ourmodel{} learns stable causal relationships that generalize across distributions. 
}

% \rln{

% \begin{table}[h]
% \scriptsize
%   \caption{Distribution Shift Performance (RMSE $\downarrow$) Comparison.}
%   \label{tab:ood_results}
%   \centering
%   \small
%   \begin{tabular}{lcc}
%     \toprule
%     Model & Exp 1 (Color) & Exp 2 (Simpson) \\
%     \midrule
%     \ourmodel{} w/o Causal & 0.204 & 0.244 \\
%     \textbf{\ourmodel{} (Ours)} & \textbf{0.198} & \textbf{0.223} \\
%     \bottomrule
%   \end{tabular}
% \end{table}

% }

% }

\subsection{Ablation Study}

\rlend{
In this section, we conduct ablation studies on the Adult and Diamonds datasets to verify the effectiveness of different components in \ourmodel{}. By comparing the full model with six ablation variants, we evaluate the contribution of each module. The results are illustrated in Figure~\ref{fig:ablation}. Specifically, we design the following ablation models: }
\vldb{
\begin{itemize}
    \item \textbf{\ourmodel{} w/o pred reg}: Removes the distance-aware regularization, treating all features equally regardless of causal distance.
    
    \item \textbf{\ourmodel{} w/o pred weight}: Eliminates the learnable weights $a_i$, treating each $h_i$ equally.
    
    \item \textbf{\ourmodel{} w/o sparse loss}: Discards the sparsity constraint on the causal mask $M$.
    
    \item \textbf{\ourmodel{} w/o acyclicity loss}: Removes the acyclicity penalty, allowing cycles to exist in the learned causal graph.
    
    \item \textbf{\ourmodel{} w/o reconstruction loss}: Omits the self-supervised reconstruction task, removing the supervision signal for learning the causal matrix.
    
    \item \textbf{\ourmodel{} w/o causal mask}: Replaces the causal-masked self-attention with a normal self-attention mechanism.
\end{itemize}
}
\rln{
As shown in Figure~\ref{fig:ablation}, removing any component from \ourmodel{} results in a noticeable performance drop, confirming the contribution of each module. 
\rlx{Specifically, removing the reconstruction loss causes the largest degradation, as the causal matrix without proper supervision may misguide the prediction. 
Removing the adaptive weights leads to significant performance drops, confirming that exploiting causal hierarchy improves performance. 
Removing the acyclicity or sparsity constraints also leads to performance degradation, demonstrating their role in ensuring valid and concise causal graph learning. Removing the prediction regularization and causal mask shows consistent drops, indicating that integrating causal structure into column-wise representation and aggregation is necessary.}
}
Overall, \ourmodel{} achieves the best results when all modules are integrated, demonstrating the effectiveness and complementarity of the proposed design.

\subsection{Effectiveness of Design Choices}
\sydesi{In this section, we evaluate the effectiveness of three key design choices in \ourmodel{} using the Adult and Diamonds datasets: the causal mask, the predictor aggregation strategy, and \syend{the reconstruction task design}.}

\begin{figure}[tb]
    \centering
    \includegraphics[width=1.0\columnwidth]{ablation_narrow_bars.png}
    \caption{Ablation study results}
    \label{fig:ablation}
    % \vspace{-5mm}
\end{figure}

% \begin{table*}[!t]
%     \caption{Comparison on Per-Epoch Training Time and Inference Time (in seconds). \\
%     ``tr'' denotes training time per epoch, and ``inf'' denotes inference time.}
%     \label{tab:efficiency}
%     \centering
%     \small
%     \setlength{\tabcolsep}{1.6mm}
%     \begin{tabular}{c*{10}{c}}
%         \toprule
%         \multirow{2}{*}{\textbf{Dataset}} &
%         \multicolumn{2}{c}{\textbf{FT-Transformer}} &
%         \multicolumn{2}{c}{\textbf{Tab-Transformer}} &
%         \multicolumn{2}{c}{\textbf{SAINT}} &
%         \multicolumn{2}{c}{\makecell{\textbf{\done{\ourmodel{}-w/o-}}\\\textbf{\done{Lazy-Update}}}} &
%         \multicolumn{2}{c}{\textbf{\ourmodel{}}} \\
%         \cmidrule(lr){2-3}\cmidrule(lr){4-5}\cmidrule(lr){6-7}\cmidrule(lr){8-9}\cmidrule(lr){10-11}
%          & \textbf{tr} & \textbf{inf} 
%          & \textbf{tr} & \textbf{inf} 
%          & \textbf{tr} & \textbf{inf} 
%          & \textbf{tr} & \textbf{inf} 
%          & \textbf{tr} & \textbf{inf} \\
%         \midrule
%         creditcard     & 25.89 & 2.98 & 11.19 & 2.18 & 21.25 & 7.39 & 25.36 & 3.53 & 17.06 & 3.45 \\
%         cardiovascular &  5.78 & 0.26 &  1.47 & 0.23 &  1.85 & 1.04 &  3.31 & 0.31 &  1.56 & 0.31 \\
%         adult          &  4.00 & 0.22 &  1.07 & 0.21 &  1.58 & 1.46 &  2.42 & 0.34 &  1.54 & 0.34 \\
%         diamonds       &  4.50 & 0.16 &  1.12 & 0.20 &  1.76 & 0.79 &  2.06 & 0.19 &  1.04 & 0.19 \\
%         elevator       &  8.65 & 0.28 &  2.08 & 0.28 &  2.54 & 0.87 &  3.76 & 0.37 &  2.05 & 0.32 \\
%         house sale     &  3.07 & 0.30 &  1.74 & 0.26 &  2.58 & 1.24 &  3.78 & 0.63 &  2.47 & 0.63 \\
%         \bottomrule
%     \end{tabular}
% \end{table*}

\subsubsection{Effectiveness of the Causal Mask}
\rlend{
For the learned causal graph \syfix{$M$ (defined in Eq~\ref{eq:masked-attn})} during attention computation, we explore three integration strategies. 
Let $A \in \mathbb{R}^{(n+1) \times (n+1)}$ denote the raw attention logits (before softmax), and $\text{Attention}$ be the normalized attention weights. }
\vldb{
\begin{itemize}
    \item \textbf{Multiplicative Masking (Ours).} Applies element-wise multiplication between the attention weights and the causal mask as defined in Equation~\ref{eq:masked-attn}.
    
    \item \textbf{Additive Masking.} Adds the causal scores to the normalized attention weights: $\tilde{\text{Attention}}_{i,j} = \text{Attention}_{i,j} + M_{i,j}$.
    
    \item \textbf{Softmax Modification.} Instead of modifying post-softmax scores, this method incorporates the causal mask into the softmax: $\text{Attention}_{i,j} = \frac{\exp(A_{i,j}) \cdot M_{i,j}}{\sum_k \exp(A_{i,k}) \cdot M_{i,k}}$.
\end{itemize}
}

\begin{table}[h]
    \centering
    \small
    \caption{Comparison of Different Mask Designs}
    \label{tab:mask_comparison}
    \begin{tabular}{lccc}
        \toprule
        Dataset & Ours & Add-Mask & Softmax-Mask \\
        \midrule
        adult   &\syne{ \textbf{0.9116}} & \syne{0.9022} & \syne{0.9040} \\
        diamonds & \syne{\textbf{0.1370}} & \syne{0.1456} & \syne{0.1445} \\
        \bottomrule
    \end{tabular}
\end{table}

% \begin{table*}[t]
%   \caption{Ablation Study on Component Designs (Adult and Diamonds)}
%   \label{tab:combined_ablation}
%   \centering
%   \begin{tabular}{l c cc cc c}
%     \toprule
%     \multirow{2}{*}{Dataset} & \multirow{2}{*}{\textbf{Ours}} & \multicolumn{2}{c}{Mask Design} & \multicolumn{2}{c}{Predictor Type} & Reconstruction \\
%     \cmidrule(lr){3-4} \cmidrule(lr){5-6} \cmidrule(lr){7-7}
%     & & Add-Mask & Softmax-Mask & Flattening & Mean Pooling & Parallel \\
%     \midrule
%     Adult & \textbf{0.9122} & 0.9062 & 0.9080 & 0.9056 & 0.9102 & 0.9085 \\
%     Diamonds & \textbf{0.0188} & 0.0195 & 0.0192 & 0.0201 & 0.0202 & 0.0191 \\
%     \bottomrule
%   \end{tabular}
% \end{table*}
% \begin{table}[t]
%   \caption{Comparison of Different Mask Designs}
%   \label{tab:mask_comparison}
%   \centering
%   \begin{tabular}{lccc}
%     \toprule
%     Dataset & Ours & Add-Mask & Softmax-Mask \\
%     \midrule
%     Adult     & \textbf{0.9122} & 0.9062 & 0.9080 \\
%     Diamonds  & \textbf{0.0188} & 0.0195 & 0.0192 \\
%     \bottomrule
%   \end{tabular}
% \end{table}

\syadd{As shown in Table~\ref{tab:mask_comparison}, multiplicative masking yields the best performance.} \rlend{This is because additive masking distorts the normalized attention distribution, while softmax modification may cause gradient vanishing when $M_{i,j}$ approaches zero. In contrast, multiplicative masking preserves the attention structure while suppressing irrelevant connections, achieving a better balance between interpretability and predictive accuracy.
}

\subsubsection{Predictor Aggregation Strategy}
Given the contextualized representations $H = \{H_1, H_2, \ldots, H_{n+1}\}$ for $n{+}1$ columns ($n$ features and the label $y$), we explore aggregation methods to obtain a unified representation $h_{\text{agg}}$ for prediction. To validate that the distance captures hierarchical feature importance, we compare the \vldb{following three strategies:

\begin{itemize}
    \item \textbf{Hierarchy-aware Aggregation (Ours).} Utilizes learnable weights $\mathbf{a}$ regularized according to the causal hierarchy, as defined in Equation~\ref{eq:distance_regularization}.
    
    \item \textbf{Flattening.} Concatenates column representations into a single vector: $h_{\text{agg}} = \text{Flatten}(H)$.
    
    \item \textbf{Mean Pooling.} Aggregates via $h_{\text{agg}} = \frac{1}{n+1}\sum_{i=1}^{n+1} H_i$, treating all columns equally.
\end{itemize}
}
\begin{table}[h]
    \centering
    \small
    \caption{Comparison of Different Predictor Aggregation Strategies}
    \label{tab:predictor_comparison}
    \begin{tabular}{lccc}
        \toprule
        Dataset & Ours & Flattening & Mean Pooling \\
        \midrule
        adult   & \syne{\textbf{0.9116}} & \syne{0.9036} & \syne{0.9062} \\
        diamonds & \syne{\textbf{0.1370}} & \syne{0.1457} & \syne{0.1441} \\
        \bottomrule
    \end{tabular}
\end{table}


\syadd{As shown in Table~\ref{tab:predictor_comparison}, hierarchy-aware aggregation achieves the most consistent improvement across tasks.} \rln{Flattening performs the worst, as concatenating all representations leads to over-parameterization and overfitting. Mean pooling shows moderate performance, but treating all columns equally ignores their hierarchical importance in the causal structure. \syadd{This confirms that features farther from the target in the causal DAG are indeed less informative than those that are causally closer.} In contrast, our hierarchy-aware aggregation dynamically adjusts each column's contribution based on its position in the learned causal hierarchy, enabling the causal structure to guide more accurate prediction.
}

\subsubsection{Reconstruction Task Design}
%%%done
\rlend{
In the causal dependency extraction module, we aim to prevent the model from using a column's own value for self-prediction. Without proper constraints, the hidden representation for reconstructing a column may contain information from itself, leading to information leakage. To address this, we explore two strategies}\syadd{: 
\begin{itemize}
    \item \textbf{\texttt{MSK} Token-based Reconstruction (Ours).}  
    In this scheme, one field is masked at a time, and its embedding is replaced by a special \texttt{MSK} token at the input layer. The model then reconstructs only the masked field.

    \item \textbf{Parallel Reconstruction.}  
    Alternatively, all fields are reconstructed simultaneously by enforcing a zero diagonal in the causal mask to block self-attention. Skip connections are removed to avoid self-prediction. While fully parallelizable, this approach may reduce training stability.
\end{itemize}
}
% \textbf{MSK Token-based Reconstruction (Ours)} masks one column at a time and reconstructs the masked column, as described in Section~\ref{sec:causal_extraction}. \textbf{Parallel Reconstruction} reconstructs all columns simultaneously by enforcing a zero diagonal in the causal mask to block self-attention; residual connections, a core design in Transformer-based models, are also removed to avoid self-prediction.

\begin{table}[h]
    \centering
    \small
    \caption{Effect of Reconstruction Task Design}
    \label{tab:reconstruction_comparison}
    \begin{tabular}{lcc}
        \toprule
        Dataset & Ours & Parallel Reconstruction \\
        \midrule
        adult   & \syne{\textbf{0.9116}} & \syne{0.9065} \\
        diamonds & \syne{\textbf{0.1370}} & \syne{0.1412} \\
        \bottomrule
    \end{tabular}
\end{table}

\syadd{As reported in Table~\ref{tab:reconstruction_comparison}, the \texttt{MSK}-based strategy consistently outperforms parallel reconstruction in both reconstruction accuracy and training stability.} \rlend{Parallel reconstruction removes residual connections to \syend{prevent} leakage, which disrupts gradient flow. In contrast, the \texttt{MSK}-based approach \syend{avoids} self-reference while preserving residual connections, leading to more stable optimization.
}

\begin{figure}[t]
  \centering
  \includegraphics[width=1.0\columnwidth]{training_time_only.png} 
  \caption{Training time comparison.}
  % \vspace{-5mm}
  \label{fig:efficiency}
\end{figure}

% \vspace{-2mm}
\subsection{Efficiency Comparison}
\syfix{
In this section, we compare the training efficiency of \ourmodel{} against all twelve baselines spanning the four categories. 
We also include \ourmodel{} w/o Lazy to validate the effect of the lazy update strategy.
Experiments are conducted on three representative datasets---Elevator (7 features), CreditCard (30 features), and MEPS (138 features)---which span small, medium, and large feature counts, and the results are illustrated in Figure~\ref{fig:efficiency}.
}
% \syfix{For clarity, let $n$ denote the number of feature columns, $d_e$ the embedding dimension, and $n_h$ the number of attention heads. \ourmodel{} shares the same training/inference complexities as Transformers ($O(n_h n^2 d_e)$ / $O(n_h n d_e)$), but the acyclicity loss adds $O(n^2)$. The lazy update strategy (Algorithm~\ref{alg:lazy_update}, every $T_c=10$ batches) amortizes this cost.}

\syfix{
% For training time, 
Among the attention-based baselines, FT-Transformer is relatively slow due to its full self-attention over all tokenized features, while Tab-Transformer is faster as it applies attention only to categorical features. 
SAINT incurs substantial training time due to its interleaved row-wise and column-wise attention. 
Orion-Bix is among the slowest methods due to the heavy bidirectional cross-attention computation. 
TabM, as an MLP-based ensemble, exhibits moderate training cost, slower than Tab-Transformer but faster than the heavier attention-based models. 
Graph-based methods exhibit mixed efficiency: ATT-Reg introduces moderate overhead from constructing inter-column adjacency matrices, while TabGNN incurs substantial overhead due to multiplex graph construction. 
Among causal-based methods, LogCause is lightweight, whereas CASTLE's cost is higher due to the matrix exponential computation in its DAG constraint; both require many epochs to converge as the complex optimization landscape slows convergence. 
LLM-CD remains relatively efficient in per-epoch training, as its LLM-based reasoning serves as a one-time preprocessing step.
}

\syfix{
Since XGBoost and AutoML do not have the concept of epochs, we thus compare \ourmodel{} with them in terms of total training time. 
On Elevator (7 features), \ourmodel{} completes in 340s compared to 160s (XGBoost) and 420s (AutoML). On CreditCard (30 features), \ourmodel{} takes 1450s compared to 1400s (XGBoost) and 1760s (AutoML). On MEPS (138 features), \ourmodel{} takes 964s compared to 990s (XGBoost) and 1170s (AutoML). The results reveal a clear trend: XGBoost is substantially faster on the small-scale dataset, but its advantage rapidly diminishes and even reverses as the feature count increases. 
This is because tree-based split search does not scale with the number of features\cite{r:xgboost}, whereas \ourmodel{}'s attention-based architecture with lazy update handles high-dimensional feature spaces more gracefully.
}

\syfix{Without the Lazy Update strategy, \ourmodel{}'s training time is high, as the acyclicity loss requires computing the matrix exponential $\operatorname{tr}(e^M)$ (Equation~\ref{eq:dag-loss}) at every update step. 
With Lazy Update, \ourmodel{} reduces training time by approximately 35\%--60\% across datasets and achieves comparable efficiency to Tab-Transformer. }
\syne{Two factors explain this improvement. First, the causal matrix $C$ stabilizes rapidly during training: after the first few epochs, the acyclicity and sparsity losses converge to near-zero values and change only marginally between consecutive mini-batches, making per-step updates of the causal dependency extraction module largely redundant.}
\syne{Second, the lazy update strategy reduces update frequency without sacrificing convergence. Unlike prediction, causal optimization tolerates sparser updates: DAG constraints and reconstruction loss ensure stability, so even infrequent updates guide the causal matrix to a valid solution. Thus, lowering frequency cuts overhead while preserving correctness and performance.}
% , as causal structure optimization benefits more from well-converged parameter states than from continuous gradient signals.
\syfix{Overall, while \ourmodel{} incorporates additional causal modeling, its computational overhead remains comparable to other baseline models.}
%%% (ok) zj2-done: delete inference
%再跑一下实验数据,training time这里多讲一些effiency多一些

% \syfix{
% \subsection{Complexity Analysis}

% Model complexity is an important consideration for practical deployment. In this section, we compare \ourmodel{} with the baselines in terms of model size. To facilitate analysis,
% we define the following notations: let \(m\) denote the number of input features, \(d_e\) the embedding dimension, \(K\) the number of attention heads, \(C\) the total number of categories across all categorical columns, and \(n_f\) the total parameters of feed-forward networks.

% All methods share a common embedding layer, whose parameter count is \(O(C d_e + m d_e)\), accounting for both categorical columns (via embedding tables) and numerical columns (via linear projection). Beyond this embedding layer, different families introduce additional parameters of different scales. For Transformer-based methods (e.g., FT-Transformer, TabTransformer, SAINT, Orion-Bix), the attention layers and feed-forward networks add \(O(K d_e^2 + n_f)\) parameters, making their total parameter count \(O(C d_e + m d_e + K d_e^2 + n_f)\). 
% For causal-based methods with MLP backbones (e.g., CASTLE, LogCause, LLM-CD), learning causal dependencies introduces a learnable causal adjacency matrix (up to \(O(m^2)\) parameters), which also subsumes the MLP parameters, leading to a total of \(O(C d_e + m d_e + m^2)\). For graph-based methods (e.g., ATT-Reg, TabGNN), modeling pairwise feature relationships directly introduces an adjacency or edge weight matrix with \(O(m^2)\) parameters, resulting in a total of \(O(C d_e + m d_e + m^2)\).   
% Tree-based models like XGBoost have parameter counts linear in the number and depth of trees, but they generally achieve lower predictive performance on the benchmarks considered in Table~\ref{tab:model_comparison}.

% \ourmodel{} shares the same embedding layer as Transformer-based baselines, with parameter count \(O(C d_e + m d_e)\). On top of this, \ourmodel{} adopts a causal-incorporated attention block, whose attention parameters are \(O(K d_e^2)\) — identical to standard Transformers. Beyond the attention mechanism, \ourmodel{} introduces a learnable causal matrix \(C \in \mathbb{R}^{(m+1)\times(m+1)}\) to capture pairwise causal dependencies, contributing \(O(m^2)\) parameters, and a hierarchy-aware prediction module that assigns learnable aggregation weights to each feature, adding \(O(m)\) parameters. These design choices enable the model to leverage causal structure and feature hierarchy for more accurate prediction.
% The total parameter count of \ourmodel{} is therefore \(O(C d_e + m d_e + K d_e^2 + m^2 + m)\). The embedding and attention terms are exactly the same as those in Transformer-based methods, and the extra \(O(m^2 + m)\) is modest, especially since the number of features \(m\) is typically small (often below \(10^3\)). Hence, \ourmodel{} achieves superior predictive performance and interpretability with negligible additional parameter cost.}

\syadd{\subsection{Hyperparameter Robustness}}
\syadd{In this section, we investigate the influence of two key hyperparameters: $\alpha$ and $\beta$ defined in Equation~\ref{eq:total_reconstruction_loss}. \syne{Specifically, $\alpha$ balances the importance between the reconstruction loss and prediction loss, and $\beta$ controls the strength of the DAG constraint.} As shown in Figure~\ref{fig:sensitive}, for $\alpha$, the performance peaks when $\alpha = 0.1$, indicating a good trade-off. A larger $\alpha$ causes the reconstruction task to dominate and harm predictive accuracy (e.g., $\alpha \geq 1$ leads to notable degradation), while a smaller $\alpha$ weakens the learning of causal structures. For $\beta$, the best performance appears around $\beta = 10^{-3}$; too large a $\beta$ (e.g., $\beta \geq 0.1$) leads to a sharp performance drop due to excessive regularization that limits model flexibility.}

\begin{figure}[htbp]
    \centering
    \includegraphics[width=0.8\columnwidth]{hyperpa.png}
    \caption{Sensitivity Analysis of Hyperparameters.}
    \label{fig:sensitive}
\end{figure}

\section{Conclusions}
\rlend{
In this paper, we propose \ourmodel{}, a framework integrating causal reasoning with attentive mechanisms for structured data prediction. \ourmodel{} learns a DAG-based causal structure via a self-supervised reconstruction task, and leverages hierarchy-aware regularization to align feature importance with causal relevance. \syne{A lazy update strategy further reduces training cost without sacrificing performance.} Experiments demonstrate that \ourmodel{} achieves superior performance and interpretability across real-world datasets.
}

% \que{In this paper, we propose \ourmodel{}, which combines a \textit{Causal Masked Attention Mechanism} and a \textit{Hierarchy-Aware \sy{Prediction module}} to jointly model causal structure and enhance structured data prediction. }\ourmodel{} alternates between a \textit{\sy{causal dependency extraction module}} for learning a DAG-based causal mask and a \textit{\sy{hierarchy-aware prediction module}} with hierarchy-aware regularization. \ourmodel{} achieves superior performance and efficiency on real-world datasets.

% \begin{acks}
% %页数：正文 ≤12页（不含参考文献和AI生成内容声明），不允许附录。ORCID：所有作者必须提供。鼓励提交代码。使用AI生成内容（文本、图、代码等）必须在致谢部分披露，说明工具、用途、具体章节。语法润色除外。摘要有字符限制，现在超了。
% \syadd{The authors confirm that no AI-generated content was used in the writing of this paper, and all text, figures, and code were independently developed and verified by the authors themselves.}
% \end{acks}

% \begin{acks}
%  This work was supported by the [...] Research Fund of [...] (Number [...]). Additional funding was provided by [...] and [...]. We also thank [...] for contributing [...].
% \end{acks}

\bibliographystyle{ACM-Reference-Format}
\bibliography{myrefs}

% --- 附录部分开始 ---
% \appendices

% \section{Methodology Details and Theoretical Proofs}
% \subsection{The Detail of Optimization with Lazy Update}
% Updating the \sy{causal dependency extraction module} can be time-consuming due to the need to compute the \sy{acyclicity loss}. To reduce the overall computational cost, we propose a \textit{lazy update strategy}. The core idea is to reduce the frequency of updating the \sy{causal dependency extraction module}. 
% %%Algorithm~\ref{alg:lazy_update} illustrates 
% \sydes{Algorithm~\ref{alg:lazy_update} shows} how the lazy update mechanism alternates training between \sy{the causal dependency extraction module and hierarchy-aware prediction module}.

% Let $C$ denote the parameters related to the causal mask, $P$ denote the predictor-related parameters (including $a_i$), and $\Theta$ represent all other model parameters. $T_s$ is the number of mini-batches \done{used} for \sy{causal dependency extraction module} updates, and $T$ is the cycle length. Specifically, for every $T$ mini-batches, the first $T_s$ are \done{\sy{employed}} to update the \sy{causal dependency extraction module}, while the remaining $T - T_s$ are \done{\sy{utilized}} to update the prediction \kdd{module}. A smaller $T_s/T$ reduces training time but may lead to lower model performance. Empirically, we set $T_s=2$ and $T=10$.

% \begin{algorithm}[h]
%   \caption{Lazy Alternating Optimization of \ourmodel{}}
%   \label{alg:lazy_update}
%   \small
%   \begin{algorithmic}[1]
%     \STATE \textbf{Input}: Training data $\mathcal{D}$, causal mask matrix $C$, \sy{hierarchy-aware prediction module} $P$, parameters $\Theta$
%     \STATE \textbf{Parameters}: Warm-up steps $T_s$, total cycle length $T$
%     \STATE \textbf{Output}: Trained \ourmodel{} model
%     \STATE Initialize iteration counter $t \leftarrow 0$
%     \FOR{each mini-batch $B \in \mathcal{D}$}
%       \IF{$t \bmod T < T_s$}
%         \STATE Randomly select a column in $B$ and replace it with \texttt{[MSK]}
%         \STATE Compute reconstruction objective $\mathcal{L}_R \leftarrow \alpha \mathcal{L}_{\text{rec}} + \beta(\mathcal{L}_{\text{DAG}} + \mathcal{L}_{\text{sparse}})$
%         \STATE Update $C$ and $\Theta$ by minimizing $\mathcal{L}_R$
%       \ELSE
%         \STATE Replace target $y$ in $B$ with \texttt{[MSK]}
%         \STATE \sydesi{Fix} causal mask $C$ and compute prediction objective $\mathcal{L}_P \leftarrow \mathcal{L}_{\text{pred}} + \mathcal{L}_{\text{dis}}$
%         \STATE Update \sy{hierarchy-aware prediction module} $P$ and $\Theta$ by minimizing $\mathcal{L}_P$
%       \ENDIF
%       \STATE $t \leftarrow t + 1$
%     \ENDFOR
%     \STATE \textbf{return} Trained \ourmodel{} model
%   \end{algorithmic}
% \end{algorithm}

% \subsection{Details of Design Choices}

% In the main text, we compared \textit{MSK Token-based Reconstruction} with \textit{Parallel Reconstruction}. While the experimental results demonstrate the superiority of the MSK-based approach, this subsection provides the theoretical justification. specifically, we mathematically derive why the \textit{Parallel Reconstruction} strategy is incompatible with residual connections, leading to a dilemma between information leakage and training instability.

% \subsubsection{Theoretical Analysis of Information Leakage in Parallel Reconstruction}

% In a standard Transformer block, the update rule for the hidden representation $H^{(l)}$ at layer $l$ typically includes a residual connection:
% \begin{equation}
%     H^{(l+1)}_i = \text{LayerNorm}\left(H^{(l)}_i + \text{Attention}(H^{(l)})_i\right)
% \end{equation}
% where $H^{(l)}_i$ is the representation of the $i$-th column.

% The \textit{Parallel Reconstruction} strategy attempts to reconstruct all columns simultaneously by masking the diagonal of the attention matrix (i.e., setting the causal mask $M_{i,i}=0$). The goal is to prevent the model from attending to the target $s_i$ itself. Under this constraint, the attention output is:
% \begin{equation}
%     \text{Attention}(H^{(l)})_i = \sum_{j=1, j \neq i}^{n} \text{AttnScore}_{i,j} \cdot V_j
% \end{equation}
% Since $M_{i,i}=0$, the term involving $V_i$ (which contains information about $s_i$) is indeed excluded from the attention aggregation.

% \textbf{The Dilemma:} To prevent this leakage in the Parallel strategy, one must remove the residual connection:
% \begin{equation}
%     H^{(l+1)}_i = \text{LayerNorm}(\text{Attention}(H^{(l)})_i)
% \end{equation}
% However, as discussed in the main text, removing residual connections significantly hinders gradient flow, making deep networks difficult to train and leading to the performance degradation observed in our ablation studies.

% \subsubsection{Justification for MSK-based Reconstruction}

% The \textit{MSK Token-based} strategy adopted in \ourmodel{} resolves this conflict fundamentally. By physically replacing the input embedding of the target with a special token:
% \begin{equation}
%     H^{(0)}_i = e_{\text{MSK}}
% \end{equation}
% The residual stream $H^{(l)}_i$ originates from $e_{\text{MSK}}$, which contains no information about the ground truth $s_i$. Therefore, even when residual connections are preserved to ensure training stability, the model is mathematically forced to aggregate information from other columns $\{s_j\}_{j \neq i}$ via the attention mechanism to minimize the reconstruction loss. This allows \ourmodel{} to benefit from both deep architecture stability and valid self-supervised causal learning.

% \section{Datasets}

% All datasets are publicly available on Kaggle or the UCI Machine Learning Repository. Data needs to be processed into CSV format, with some datasets requiring special processing. The specific download links and preprocessing details for each dataset are as follows:

% \subsection{Adult Dataset}
% \textbf{Source}: \url{https://archive.ics.uci.edu/dataset/2/adult} \\
% \textbf{Task}: Binary classification (income prediction) \\
% \textbf{Preprocessing}: None

% \subsection{\syx{Cardiovascular} Dataset}
% \textbf{Source}:\url{https://www.kaggle.com/datasets/sulianova/cardiovascular-disease-dataset} \\
% \textbf{Task}: \syx{Binary classification} (cardiotocography classification) \\
% \textbf{Preprocessing}: Four features: height, weight, ap\_hi, ap\_lo are discretized into categorical variables with intervals of 10

% \subsection{CreditCard Dataset}
% \textbf{Source}:\url{https://www.kaggle.com/datasets/arockiaselciaa/creditcardcsv} \\
% \textbf{Task}: Binary classification (credit card \syx{fraudulent transactions detection}) \\
% \textbf{Preprocessing}: None

% \subsection{Diamonds Dataset}
% \textbf{Source}:\url{https://www.kaggle.com/datasets/shivam2503/diamonds} \\
% \textbf{Task}: Regression (diamond price prediction) \\
% \textbf{Preprocessing}: The price column is moved to the last column as the prediction target

% \subsection{Elevator Dataset}
% \textbf{Source}:\url{https://www.kaggle.com/datasets/shivamb/elevator-predictive-maintenance-dataset} \\
% \textbf{Task}: Regression (elevator behavior prediction) \\
% \textbf{Preprocessing}: The vibration column is moved to the last column as the prediction target

% \subsection{Housesale Dataset}
% \textbf{Source}:\url{https://www.kaggle.com/datasets/ruchi798/housing-prices-in-metropolitan-areas-of-india} \\
% \textbf{Task}: Regression (house price prediction) \\
% \textbf{Preprocessing}: All 6 CSV files in the dataset are concatenated together, and the data is shuffled with random seed 42

% \subsection{Crime Dataset}
% \textbf{Source}: \url{https://archive.ics.uci.edu/dataset/183/communities+and+crime} \\
% \textbf{Task}: Regression (violent crime rate prediction) \\
% \textbf{Preprocessing}: Non-predictive features (e.g., community name and state) are removed, and missing values are handled via mean imputation

% \subsection{MEPS Dataset}
% \textbf{Source}: \url{https://meps.ahrq.gov/mepsweb/} \\
% \textbf{Task}: \syx{Regression} (healthcare utilization prediction) \\
% \textbf{Preprocessing}: \syx{Using the scripts from  \url{https://github.com/yromano/cqr/tree/master/get_meps_data}, we merged MEPS data from Panel 19 (2015, $n=15,785$), Panel 20 (2015, $n=17,541$), and Panel 21 (2016, $n=15,656$)}

% \section{Implementation Details}

% All methods are implemented using PyTorch 2.2.2 with CUDA 12.2, and the experiments are conducted on a server equipped with an Intel(R) Xeon(R) w7-3455 CPU and two NVIDIA RTX A6000 GPUs.In all experiments, the batch size is set to 128. The optimizer employed is Adam with a weight decay of $1 \times 10^{-5}$. Other hyperparameters are selected based on the best-performing configuration for each model on each dataset.

% We employ AUROC and RMSE to evaluate the model performance on classification and regression tasks, respectively. A higher AUROC indicates better classification performance, while a lower RMSE implies better regression performance. 

% Additionally, the dropout rate is set to 0.1. For models based on multi-head self-attention, the number of heads is set to the hidden size divided by 8. The hyperparameters are specifically selected via grid search on the validation dataset, and the selected values are then adopted to report the test performance.

% In the synthetic data experiments, the datasets are generated based on the NPSEM framework. The graph structures are set to 5-node DAGs with 6 edges and 10-node DAGs with 15 edges. The noise terms are sampled from a Gaussian distribution, and the sample size is determined to ensure statistical significance.

% For the distribution shift evaluation, the dataset splitting protocols are set to strictly simulate out-of-distribution scenarios. In the feature-level shift, the training set is composed of 90\% of specific categories, while the test set is dominated by the remaining ones. The splitting configuration is selected to assess robustness against spurious correlations.

% For \textbf{\ourmodel{}}, the grid for hidden size is [16, 32, 64, 128], and 32 is selected as optimal. The batch size is chosen from [64, 128, 256, 512], with 128 being best. The learning rate is selected from [1e-5, 1e-4, 1e-3, 1e-2], and 1e-3 is found optimal. The number of encoder layers is chosen from [1, 2, 3, 4, 5, 6], with 2 selected. 
% % The number of epochs is set to 20.

% For \textbf{FT-Transformer} \citep{r:ft-transformer}, the hidden size is selected from [16, 32, 64, 128], and 64 is chosen. The batch size is selected from [64, 128, 256, 512], with 128 being best. The learning rate is chosen from [1e-5, 1e-4, 1e-3, 1e-2], and 1e-3 is selected. The number of encoder layers is chosen from [1, 2, 3, 4, 5, 6], and 2 is selected. % The number of epochs is 30.

% For \textbf{Tab-Transformer} \citep{r:tab}, the hidden size is selected from [16, 32, 64, 128], with 32 being optimal. The batch size is chosen from [64, 128, 256, 512], and 128 is used. The learning rate is selected from [1e-5, 1e-4, 1e-3, 1e-2], and 1e-3 is best. The encoder layers are chosen from [1, 2, 3, 4, 5, 6], and 2 is selected.% The number of epochs is 30.

% For \textbf{SAINT} \citep{r:SAINT}, the hidden size is selected from [16, 32, 64, 128], and 64 is best. The batch size is selected from [64, 128, 256, 512], with 128 chosen. The learning rate is selected from [1e-5, 1e-4, 1e-3, 1e-2], and 1e-4 is found optimal. The number of encoder layers is selected from [2, 4, 6, 8], and 2 is chosen (as each SAINT block contains both row and column attention).% The number of epochs is 50.

% For \textbf{CASTLE} \citep{r:CASTLE}, the hidden size is selected from [16, 32, 64, 128], and 32 is best. The batch size is chosen from [64, 128, 256, 512], and 64 is optimal. The learning rate is selected from [1e-5, 1e-4, 1e-3, 1e-2], and 1e-4 is found optimal. The number of MLP layers is selected from [1, 2, 3, 4], and 2 is chosen. The special hyperparameters $\alpha$ and $\beta$ are set according to the paper's recommendation as 1 and 5, respectively.% The number of epochs is 50.

% For \textbf{LogCause} \citep{r:linear}, the hidden size is selected from [16, 32, 64, 128], and 32 is used. The batch size is selected from [64, 128, 256, 512], with 64 chosen. The learning rate is selected from [1e-5, 1e-4, 1e-3, 1e-2], and 1e-5 is optimal.% The number of epochs is 200.

% For \textbf{XGBoost} \citep{r:xgboost}, the max depth is selected from [3, 6, 9], and 6 is found optimal. The learning rate is chosen from [1e-2, 5e-2, 1e-1], and 0.1 is best. The number of trees is selected from [100, 500, 1000], and 500 is chosen. Both the subsample and colsample by tree are chosen from [0.6, 0.8, 1.0], and 0.8 is optimal for both.

% For \textbf{AutoML} \citep{r:automl}, Seeds are selected from [42, 142, 127, 723, 2025], and a dedicated validation set is chosen as tuning data.

% For \textbf{ATT-Reg} \citep{r:attreg}, the embedding dimension is selected from [8, 16, 32, 64], and 32 is best. The hidden size is chosen from [16, 32, 64, 128], and 32 is found optimal. The learning rate is selected from [1e-4, 5e-4, 1e-3, 5e-3], and 1e-3 is best. The batch size is chosen from [64, 128, 256, 512], and 128 is used. The weight decay is selected from [1e-7, 1e-6, 1e-5, 1e-4, 1e-3], and 1e-5 is chosen. The regularization coefficient $\lambda$ is selected from [1e-7, 1e-6, 1e-5, 1e-4, 1e-3, 1e-2], and 1e-4 is found optimal.% The number of epochs is 100.

% For \textbf{TabGNN} \citep{r:gnnbased}, the hidden size is selected from [64, 128, 256], and 128 is best. The number of GNN message passing layers is chosen from [1, 2, 3], and 2 is optimal. The learning rate is selected from [5e-5, 1e-4, 1e-3], and 1e-3 is found optimal. The batch size is chosen from [32, 256, 1024], and 256 is used. Graph construction and readout are set as categorical feature sharing and GlobalAttentionPooling, respectively.% The number of epochs is 100.

% \begin{figure}[t]
%     \centering
%     \includegraphics[width=0.8\columnwidth]{sensitive2.png}
%     \caption{Sensitivity analysis of hyperparameters $\alpha$ and $\beta$.}
%     \label{fig:sensitive}
% \end{figure}

% \section{Additional Experimental Results}

% \subsection{Sensitivity Analysis of Hyperparameters}
% In this section, we investigate the influence of two key hyperparameters: $\alpha$ and $\beta$ defined in Equation ~\ref{eq:total_reconstruction_loss}. Specifically, $\alpha$ balances the importance between the reconstruction loss and prediction loss, and $\beta$ controls the strength of the DAG constraint. 

% As shown in Figure~\ref{fig:sensitive}, for $\alpha$, the performance peaks when $\alpha = 0.1$, indicating a good trade-off. A larger $\alpha$ causes the reconstruction task to dominate and harm predictive accuracy, while a smaller $\alpha$ weakens the learning of causal structures. For $\beta$, the best performance appears around $\beta = 10^{-3}$, while too large a $\beta$ (e.g., $10^{-1}$) leads to a sharp performance drop due to excessive regularization that limits model flexibility.

%\clearpage

\end{document}
\endinput